% Chapter 2, Section 1 _Linear Algebra_ Jim Hefferon
%  http://joshua.smcvt.edu/linearalgebra
%  2001-Jun-10
\chapter{向量空间}
第一章结束时，我们已经对高斯消元法如何求解线性方程组有了相当好的理解。
它系统地取各行的线性组合。
这里我们将转向对线性组合的一般研究。

我们需要一个舞台。
第一章中，我们有时组合来自 $\Re^2$ 的向量，
有时组合来自 $\Re^3$ 的向量，
还有时组合来自更高维空间的向量。
于是我们的第一反应可能是在 $\Re^n$ 中工作，
而不指明 $n$ 取何值。
这样做有一个好处：
所有结果可以同时对 $\Re^2$、对 $\Re^3$
以及对许多其他空间成立。

但是，如果让结果同时适用于许多空间是有利的，
那么只局限于 $\Re^n$ 这类空间就太受束缚了。
我们希望我们的结果也能适用于行向量的组合，
如第一章最后一节那样。
我们甚至已经见过一些空间，
它们并非就是同等尺寸的列向量或行向量全体的集合。
例如，我们见过某个齐次方程组的解集，
它是 $\Re^3$ 内的一个平面。
这个集合是一个封闭系统，因为
这些解的线性组合仍是解。
但它并不包含所有高为3的列向量，
只包含其中的一部分。

我们希望关于线性组合的结果，
在一切线性组合有意义的场合都适用。
我们将把任何这样的集合称为\definend{向量空间}。
我们的结果将不再表述为
“每当我们有一个可以合理地取线性组合的集合\ldots”，
而是表述为“在任何向量空间中\ldots”

这样一句话同时描述了许多空间中的情形。
为了理解从研究单个空间进到研究一类空间的好处，
请看这个类比。
设想政府一次只给一个人立法：
“Leslie Jones 不得乱穿马路。”
那会很糟糕；
陈述的优点在于其经济性，即同时适用于许多情形。
又假设他们说：“Kim Ke 路过事故现场时必须停车。”
对比一下：“任何医生路过事故现场时必须停车。”
更一般的陈述，在某些方面，反而更清晰。













\section{向量空间的定义}
\index{vector space|(}
我们将研究带有两个运算的结构：
一个加法与一个标量乘法，
它们须满足一些简单的条件。
这些条件我们稍后还会进一步思考，
但初读时请注意它们是多么合理。
例如，任何可以称作加法的运算
（如列向量加法、行向量加法或
实数加法）必定满足下面的条件 (1) 至~(5)。



\vspace*{1ex}
\subsection{定义与例题}

\begin{definition}
\label{def:VecSpace}
%<*df:VectorSpace>
一个\definend{向量空间}\index{vector space!definition}
（在 \( \Re \) 上）由一个集合 \( V \) 连同
两个运算
“+”\index{vector!sum}\index{sum!vector}\index{addition of vectors} 
与“\( \cdot \)”\index{scalar multiple!vector}\index{vector!scalar multiple} 
组成，它们满足以下条件：
对所有向量 \( \vec{v},\vec{w},\vec{u}\in V \)
及所有\definend{标量}\index{scalar}
\( r,s\in\Re \)：
% \begin{tfae}
\begin{compactenum} 
\item 集合 $V$ 对
  向量加法封闭，也就是说，
  \( \vec{v}+\vec{w}\in V \)
\item 向量加法是交换的，
  \( \vec{v}+\vec{w}=\vec{w}+\vec{v} \) 
\item 向量加法是结合的，
  \( (\vec{v}+\vec{w})+\vec{u}=\vec{v}+(\vec{w}+\vec{u}) \)
\item 存在一个\definend{零向量}\index{zero vector}
    \( \zero\in V \)，使得
    \( \vec{v}+\zero=\vec{v}\, \) 对所有 \( \vec{v}\in V\/ \) 成立
\item 每个 \( \vec{v}\in V \) 都有一个
    \definend{加法逆元}\index{additive inverse}\index{inverse!additive}
    \( \vec{w}\in V \)，使得 \( \vec{w}+\vec{v}=\zero \)
\item 集合 $V$ 对
    标量乘法封闭，也就是说，
   \( r\cdot\vec{v}\in V \)
\item 标量乘法对标量加法可分配，
 \( (r+s)\cdot\vec{v}=r\cdot\vec{v}+s\cdot\vec{v} \)
\item 标量乘法对向量加法可分配，
  \( r\cdot(\vec{v}+\vec{w})=r\cdot\vec{v}+r\cdot\vec{w} \)
\item 标量的普通乘法与
  标量乘法满足结合律，\( (rs)\cdot\vec{v} =r\cdot(s\cdot\vec{v}) \)
\item 用标量~$1$ 作乘法是
  恒等运算，\( 1\cdot\vec{v}=\vec{v} \)。
\end{compactenum}
% \end{tfae}  
%</df:VectorSpace>
\end{definition}

\vspace{-4.5ex}  % why is the extra space there?  Bug in shading code?
\begin{remark}
这个定义涉及两种加法与两种乘法，
因此乍看之下可能令人困惑。
例如，在条件~(7) 中，
左边的“$+$”是两个实数的加法，
而右边的“$+$”是
\( V\/ \) 中两个向量的加法。
这些表达式因上下文而无歧义；例如，
\( r \) 与 \( s \)
是实数，所以“\( r+s \)”只能表示实数加法。
同样，条目~(9) 左边的“$rs$”是普通的实数
乘法，而右边的“$s\cdot\vec{v}$”是
为这个向量空间定义的标量乘法。
\end{remark}

\vspace{-0.5ex}  % why is the extra space there?  Bug in shading code?
要理解这个定义，最好的方法是逐个考察下面的例子，
并对每个例子验证全部十条条件。
第一个例子包含这一验证，写得十分详细。
请把它作为其余各例的范本。
特别重要的是\definend{封闭性}\index{vector space!closure}
条件，即 (1) 与~(6)。
它们规定加法与标量乘法运算总是有意义的\Dash 即它们对每一对向量、
每一个标量与向量都有定义，
并且运算的结果是该集合的成员。


\begin{example}  \label{PlaneThruOriginSubsp}
\( \Re^2 \) 的这个子集是过原点的一条直线。
\begin{equation*}
  L=\set{ \colvec{x \\ y}  \suchthat y=3x}
\end{equation*}
我们将验证，在“+”与“\(\cdot\)”的通常意义下
它构成向量空间。
\begin{equation*}
  \colvec{x_1 \\ y_1}
  +
  \colvec{x_2 \\ y_2}
  =
  \colvec{x_1+x_2 \\ y_1+y_2}
  \qquad
  r\cdot
  \colvec{x \\ y}
  =
  \colvec{rx \\ ry}
\end{equation*}
这些运算
就是通常的运算，只是在它的子集 $L$ 上重新使用。
我们说 \( L \) 从 \( \Re^2 \)
\definend{继承\/}\index{inherited operations} 这些运算。

我们将验证全部十条条件。
与加法有关的条件共有五条。
对于条件~(1)，即对加法封闭，
假设我们从直线~$L$ 上取两个向量，
\begin{equation*}
  \vec{v}_1=\colvec{x_1 \\ y_1}
  \quad
  \vec{v}_2=\colvec{x_2 \\ y_2}
\end{equation*}
于是它们满足限制
$y_1=3x_1$ 与~$y_2=3x_2$。
它们的和
\begin{equation*}
  \vec{v}_1+\vec{v}_2
  =\colvec{x_1+x_2 \\ y_1+y_2}
\end{equation*}
也在直线 $L$ 上，
因为其第二个分量是第一个分量的三倍
$y_1+y_2=3(x_1+x_2)$，这由
$\vec{v}_1$ 与~$\vec{v}_2$ 所受的限制可得。
对于~(2)，即向量加法可交换，
只需比较
\begin{equation*}
  \vec{v}_1+\vec{v}_2
  =\colvec{x_1+x_2 \\ y_1+y_2}
  \quad
  \vec{v}_2+\vec{v}_1
  =\colvec{x_2+x_1 \\ y_2+y_1}
\end{equation*}
并注意它们相等，因为它们的分量是实数而
实数可交换。
（向量满足位于该直线上的限制与本条件无关；
它们可交换只是因为平面中的所有向量都可交换。）
条件~(3)，即向量加法的结合律，是类似的。
\begin{align*}
  (\colvec{x_1 \\ y_1}
  +\colvec{x_2 \\ y_2})
  +\colvec{x_3 \\ y_3}
  &=\colvec{(x_1+x_2)+x_3 \\ (y_1+y_2)+y_3}  \\
  &=\colvec{x_1+(x_2+x_3) \\ y_1+(y_2+y_3)}  \\
  &=\colvec{x_1 \\ y_1}
  +(\colvec{x_2 \\ y_2}
  +\colvec{x_3 \\ y_3})
\end{align*}
对于第四条条件，我们必须给出一个充当
零元的向量。
全部分量为零的向量即可。
\begin{equation*}
  \colvec{x \\ y}
  +\colvec{0 \\ 0}
  =\colvec{x \\ y}
\end{equation*}
注意 $\zero\in L$，因为其第二个分量是第一个分量的三倍。
对于~(5)，即给任意~$\vec{v}\in L$ 都能
产生一个加法逆元，我们有
\begin{equation*}
  \colvec{-x \\ -y}
  +\colvec{x \\ y}
  =\colvec{0 \\ 0}
\end{equation*}
所以向量 $-\vec{v}$ 就是想要的逆元。
与前一条条件一样，这里注意若 $\vec{v}\in L$，即
$y=3x$，则 $-\vec{v}\in L$ 也成立，因为 $-y=3(-x)$。

与标量乘法有关的五条条件的验证是类似的。
对于~(6)，即对标量乘法封闭，
假设 $r\in\R$ 且 $\vec{v}\in L$，即
\begin{equation*}
  \vec{v}=\colvec{x \\ y}
\end{equation*}
满足~$y=3x$。
那么
\begin{equation*}
  r\cdot\vec{v}=r\cdot\colvec{x \\ y}
  =\colvec{rx \\ ry}
\end{equation*}
也在 $L$ 中：~关系式 $ry=3\cdot rx$ 成立，因为
$y=3x$。
接下来，这验证了~(7)。
\begin{equation*}
  (r+s)\cdot\colvec{x \\ y}
  =\colvec{(r+s)x \\ (r+s)y}
  =\colvec{rx+sx \\ ry+sy}
  =r\cdot\colvec{x \\ y}+s\cdot\colvec{x \\ y}
\end{equation*}
对于~(8)，我们有：
\begin{equation*}
  r\cdot(\colvec{x_1 \\ y_1}+\colvec{x_2 \\ y_2})
  =\colvec{r(x_1+x_2) \\ r(y_1+y_2)}
  =\colvec{rx_1+rx_2 \\ ry_1+ry_2}
  =r\cdot\colvec{x_1 \\ y_1}+r\cdot\colvec{x_2 \\ y_2}
\end{equation*}
第九条
\begin{equation*}
  (rs)\cdot\colvec{x \\ y}
  =\colvec{(rs)x \\ (rs)y}
  =\colvec{r(sx) \\ r(sy)}
  =r\cdot(s\cdot\colvec{x \\ y})
\end{equation*}
与第十条条件也同样直接。
\begin{equation*}
  1\cdot\colvec{x \\ y}
  =\colvec{1x \\ 1y}
  =\colvec{x \\ y}
\end{equation*}
\end{example}

\begin{example}  \label{ex:RealVecSpaces}
整个平面，即集合
\( \Re^2 \)，是一个向量空间，其中运算“\( + \)”与“\( \cdot \)”
取通常的意义。
\begin{equation*}
  \colvec{x_1 \\ y_1}
  +
  \colvec{x_2 \\ y_2}
  =
  \colvec{x_1+x_2 \\ y_1+y_2}
  \qquad
  r\cdot
  \colvec{x \\ y}
  =
  \colvec{rx \\ ry}
\end{equation*}
我们只验证其中的两条条件，即封闭性条件。

对于 (1)，注意向量和
\begin{equation*}
  \colvec{x_1 \\ y_1}
  +\colvec{x_2 \\ y_2}
  =\colvec{x_1+x_2 \\ y_1+y_2}
\end{equation*}
的结果是一个带有两个实分量的列阵列，所以是
平面~\( \Re^2 \) 的成员。
与前面的例题相比，这里对向量的
第一个与第二个分量没有任何限制。

条件 (6) 类似。
向量
\begin{equation*}
  r\cdot\colvec{x \\ y}
  =\colvec{rx \\ ry}
\end{equation*}
有两个实分量，所以是~\( \Re^2 \) 的成员。
\end{example}

类似地，
每个 \( \Re^n \) 都是向量空间，带有通常的向量加法
与标量乘法运算。
（在 \( \Re^1 \) 中，我们通常不把成员写成
列向量，也就是说，我们通常不写“\( (\pi) \)”。
而是直接写“\( \pi \)”。）

\begin{example} \label{ex:ColsIntEntNotVS}
\nearbyexample{PlaneThruOriginSubsp} 给出了 $\Re^2$ 的一个
构成向量空间的子集。
% \nearbyexample{ex:RealVecSpaces} shows that the entire
% set of two-tall vectors with real entries is also a vector space. 
作为对比，考虑
由分量为整数的
高为2的列组成的集合，
运算仍是同样的按分量的
加法与标量乘法。
它是 $\Re^2$ 的一个子集，但它不是向量空间：
它对标量乘法不封闭，
也就是说，它不满足条件~(6)。
例如，下面的左边是一个具有整数分量的向量，以及一个标量。
\begin{equation*}
  0.5
  \cdot
  \colvec[r]{4 \\ 3}
  =
  \colvec[r]{2 \\ 1.5}
\end{equation*}
右边是一个不属于该集合的列向量，因为
它的分量不全是整数。
\end{example}

\begin{example}  \label{ex:TrivSbspReFour}
单元素集
\begin{equation*}
  \set{ \colvec[r]{0 \\ 0 \\ 0 \\ 0} }
\end{equation*}
在下述运算下构成向量空间
\begin{equation*}
  \colvec[r]{0 \\ 0 \\ 0 \\ 0}
  +
  \colvec[r]{0 \\ 0 \\ 0 \\ 0}
  =
  \colvec[r]{0 \\ 0 \\ 0 \\ 0}
  \qquad
  r\cdot
  \colvec[r]{0 \\ 0 \\ 0 \\ 0}
  =
  \colvec[r]{0 \\ 0 \\ 0 \\ 0}
\end{equation*}
这些运算继承自 \( \Re^4 \)。
\end{example}

向量空间必须至少有一个元素，即它的零向量。
因此，单元素向量空间是可能的最小情形。

\begin{definition} \label{df:TrivialVectorSpace}
%<*df:TrivialVectorSpace>
单元素向量空间是\definend{平凡的}\index{vector space!trivial}%
\index{trivial space}
空间。
%</df:TrivialVectorSpace>
\end{definition}

迄今的例子涉及的都是带有通常运算的列向量的集合。
但向量空间不必是列向量的集合，甚至不必是行
向量的集合。
下面是其他一些类型的向量空间。
“向量空间”一词并不意味着“实数列的集合”。
它的含义更像是
“任何线性组合都有意义的集合”。

\begin{example} \label{ex:PolySpaceThree}
考虑
\( \polyspace_3=\set{a_0+a_1x+a_2x^2+a_3x^3\suchthat a_0,\ldots,a_3\in\Re} \)，
即次数不超过三的多项式之集
（在本书中，我们把常数多项式，包括零多项式，
取为零次）。
在以下运算下它是向量空间
\begin{multline*}
   (a_0+a_1x+a_2x^2+a_3x^3)+(b_0+b_1x+b_2x^2+b_3x^3)  \\
     =(a_0+b_0)+(a_1+b_1)x+(a_2+b_2)x^2+(a_3+b_3)x^3
\end{multline*}
与
\begin{equation*}
   r\cdot(a_0+a_1x+a_2x^2+a_3x^3)=
     (ra_0)+(ra_1)x+(ra_2)x^2+(ra_3)x^3
\end{equation*}
（验证很容易）。
这个向量空间值得关注，因为
这些是高中代数中熟悉的多项式运算。
例如，
$
  3\cdot(1-2x+3x^2-4x^3)-2\cdot(2-3x+x^2-(1/2)x^3)=-1+7x^2-11x^3$。

虽然这个空间不是任何 \( \Re^n \) 的子集，
但在某种意义下我们可以把 $\polyspace_3$ 看作与
\( \Re^4 \) “相同”。
如果我们按如下方式将这两个空间的元素对应起来
\begin{equation*}
  a_0+a_1x+a_2x^2+a_3x^3
  \quad\text{对应于}\quad
  \colvec{a_0 \\ a_1 \\ a_2 \\ a_3}
\end{equation*}
那么运算也相互对应。
下面是一个对应加法的例子。
\begin{equation*}
  \mbox{\begin{tabular}{lr}
       &\(1-2x+0x^2+1x^3\) \\
     + &\(2+3x+7x^2-4x^3\) \\ \hline
       &\(3+1x+7x^2-3x^3\)
  \end{tabular}  }
  \quad\text{对应于}\quad
  \colvec[r]{1 \\ -2 \\ 0 \\ 1}
  +
  \colvec[r]{2 \\ 3 \\ 7 \\ -4}
  =
  \colvec[r]{3 \\ 1 \\ 7 \\ -3}
\end{equation*}
我们视为“相同”的事物相加得到“相同”的和。
第三章将把这种向量空间对应的思想精确化。
现在我们只把它当作一种直观。
\end{example}

一般地，我们用
\( \polyspace_n \) 表示
次数不超过~$n$ 的多项式组成的向量空间
\( \set{a_0+a_1x+a_2x^2+\cdots+a_nx^n\suchthat a_0,\ldots,a_n\in\Re} \)，
其运算为通常的多项式加法与标量
乘法。\index{vector space!of polynomials}
我们会经常使用这些空间作为例子。

\begin{example} \label{ex:MatSpaceTwoByTwo}
\( \nbym{2}{2} \) 的实元素矩阵之集 \( \matspace_{\nbym{2}{2}} \)
在自然的逐元素
运算下构成向量空间。
\begin{equation*}
  \begin{mat}
    a  &b \\
    c  &d
  \end{mat}
  +
  \begin{mat}
    w  &x \\
    y  &z
  \end{mat}
  =
  \begin{mat}
    a+w  &b+x \\
    c+y  &d+z
  \end{mat}
  \qquad
  r\cdot
  \begin{mat}
    a  &b \\
    c  &d
  \end{mat}
  =
  \begin{mat}
    ra  &rb \\
    rc  &rd
  \end{mat}
\end{equation*}
与前面的例题一样，我们可以认为这个空间
与 \( \Re^4 \) 是“相同的”。
\end{example}

我们用
\( \matspace_{\nbym{n}{m}} \) 表示
$\nbym{n}{m}$ 矩阵
在自然的矩阵加法与标量
乘法运算下的向量空间。\index{vector space!of matrices}
与多项式空间一样，
我们会经常用它们作为例子。

\begin{example}   \label{ex:FcnsNToRIsVecSp}
全体以一个自然数为自变量的实值函数之集
\( \set{f\suchthat \map{f}{\N}{\Re} } \)
在以下运算下构成向量空间
\begin{equation*}
  (f_1+f_2)\,(n)=f_1(n)+f_2(n)
  \qquad
  (r\cdot f)\,(n)=r\,f(n)
\end{equation*}
于是，例如若 \( f_1(n)=n^2+2\sin(n) \) 且
\( f_2(n)=-\sin(n)+0.5 \)，
则 \( (f_1+2f_2)\,(n)=n^2+1 \)。

我们可以把这个空间看作
\nearbyexample{ex:RealVecSpaces} 的推广\Dash 这些函数不是
高为$2$的向量，而像是无限高的
向量。
\begin{equation*}
  \text{
    \begin{tabular}{c|c}
      \( n \)      &\( f(n)=n^2+1 \)  \\ \hline
      \( 0      \) &\( 1      \)    \\
      \( 1      \) &\( 2      \)    \\
      \( 2      \) &\( 5      \)    \\
      $3$          &$10$            \\
      \( \vdots \) &\( \vdots \)    
    \end{tabular} }
    \quad\text{对应于}\quad
    \colvec[r]{
           1           \\
           2           \\
           5           \\
           10          \\
           \vdotswithin{10}      }
\end{equation*}
加法与标量乘法都是按分量的，
如\nearbyexample{ex:RealVecSpaces} 中那样。
（我们可以把“无限高”形式化：说它表示一个无穷
序列，或者说它表示一个从 $\N$ 到 $\Re$ 的函数。）
\end{example}

\begin{example} \label{ex:PolysOfAllFiniteDegrees}
实系数多项式之集
\begin{equation*}
 \set{ a_0+a_1x+\cdots+a_nx^n\suchthat n\in\N
    \text{ 且 } a_0,\ldots,a_n\in\Re} 
\end{equation*}
在赋予自然的“$+$”
\begin{multline*}
  (a_0+a_1x+\cdots+a_nx^n)+(b_0+b_1x+\cdots+b_nx^n)  \\
     =(a_0+b_0)+(a_1+b_1)x+\cdots +(a_n+b_n)x^n
\end{multline*}
与“$\cdot$”之后
\begin{equation*}
  r\cdot (a_0+a_1x+\ldots a_nx^n)
   =
  (ra_0)+(ra_1)x+\ldots (ra_n)x^n
\end{equation*}
便构成向量空间。
这个空间不同于
\nearbyexample{ex:PolySpaceThree} 的空间 $\polyspace_3$。
这个空间包含的不只是三次多项式，
还有三十次的与三百次的
多项式。
当然，每个多项式的次数都是有限的，
但该集合对其所有成员的次数没有一个统一的界。

与前面的例题一样，我们可以
用无穷元组来思考这个例子。
例如，我们可以认为 \( 1+3x+5x^2 \) 对应于
\( (1,3,5,0,0,\ldots) \)。
然而，这个空间不同于
\nearbyexample{ex:FcnsNToRIsVecSp} 中的空间。
这里，集合的每个成员都有有限的次数，也就是说，
在这个对应下，本空间中没有元素
与 \( (1,2,5,10,\,\ldots\,) \) 匹配。
这个空间中的向量对应于
以零结尾的无穷元组。
\end{example}

\begin{example}  \label{ex:RealValuedFcns}
全体以一个实数为自变量的实值函数之集
\( \set{f\suchthat \map{f}{\Re}{\Re} } \)
在这些运算下构成向量空间。
\begin{equation*}
  (f_1+f_2)\,(x)=f_1(x)+f_2(x)
  \qquad
  (r\cdot f)\,(x)=r\,f(x)
\end{equation*}
它与\nearbyexample{ex:FcnsNToRIsVecSp} 的差别在于
函数的定义域。
\end{example}

\begin{example}\label{ex:ACos+BSin}
实变量 \( \theta \) 的实值函数之集
\( F=\{ a\cos\theta+b\sin\theta \suchthat a,b\in\Re\} \)
在以下运算下构成向量空间
%\nearbyexample{ex:RealValuedFcns}:
\begin{equation*}
  (a_1\cos\theta+b_1\sin\theta)+(a_2\cos\theta+b_2\sin\theta)
    =(a_1+a_2)\cos\theta+(b_1+b_2)\sin\theta
\end{equation*}
与
\begin{equation*}
  r\cdot (a\cos\theta+b\sin\theta)
   =(ra)\cos\theta+(rb)\sin\theta
\end{equation*}
这些运算继承自前面例题中的空间。
（我们可以把 \( F \) 看作与 \( \Re^2 \) “相同”，
因为 $a\cos\theta+b\sin\theta$ 对应于
分量为 $a$ 和 $b$ 的向量。）
\end{example}

\begin{example}  \label{ex:SpaceSatisfyingDiffQ}
集合
\begin{equation*}
  \set{\map{f}{\Re}{\Re}\suchthat \dfrac{d^2f}{dx^2}+f=0}
\end{equation*}
在如今已很自然的解释下构成向量空间。
\begin{equation*}
  (f+g)\,(x)=f(x)+g(x)
  \qquad
  (r\cdot f)\,(x)=r\,f(x)
\end{equation*}
特别地，注意初等微积分给出
\begin{equation*}
   \frac{d^2(f+g)}{dx^2}+(f+g)
   =(\frac{d^2f}{dx^2}+f)+(\frac{d^2g}{dx^2}+g)
\end{equation*}
以及
\begin{equation*}
   \frac{d^2(rf)}{dx^2}+(rf)
   =r(\frac{d^2 f}{dx^2}+f)
\end{equation*}
所以该空间对加法与标量乘法封闭。
结果这个空间就等于前面例题中的空间\Dash 满足
此微分方程的函数形如
$a\cos\theta+b\sin\theta$\Dash 但这一描述提示可将其推广到其他
微分方程的解集。
\end{example}

\begin{example} \label{ex:HomoSlnMakesVS}
\( n \) 个变量的齐次线性方程组的解集，
在从 \( \Re^n \) 继承的运算下构成向量空间。
例如，为验证对加法封闭，
考虑该方程组中一个典型的方程
$c_1x_1+\cdots+c_nx_n=0$，并假设这两个向量
\begin{equation*}
   \vec{v}=\colvec{v_1 \\ \vdotswithin{v_1} \\ v_n}
   \qquad
   \vec{w}=\colvec{w_1 \\ \vdotswithin{w_1} \\ w_n}
\end{equation*}
都满足该方程。
那么它们的和
\( \vec{v}+\vec{w}\/ \) 也满足该方程：
\(
  c_1(v_1+w_1)+\cdots+c_n(v_n+w_n)
  =(c_1v_1+\cdots+c_nv_n)+(c_1w_1+\cdots+c_nw_n)
  =0
\)。
其余向量空间条件的验证同样常规。
\end{example}

在标量与向量之间，我们常省略
乘号“\( \cdot \)”。
我们凭上下文区分
\( c_1v_1 \) 中的乘法与 \( r\vec{v}\, \) 中的乘法，因为若两个
相乘的都是实数，则必是
实数与实数的乘法；而若其中一个是向量，则必是
标量与向量的乘法。

\nearbyexample{ex:HomoSlnMakesVS} 让我们回到了出发点，因为它是我们的
动机性例子之一。
现在，对满足向量空间
定义的各类结构有了些许感受之后，
我们可以反思这个定义。
例如，为什么在定义中规定条件
\( 1\cdot\vec{v}=\vec{v} \)，
而不规定条件 \( 0\cdot\vec{v}=\zero \)？

一个回答是：这只是一个定义\Dash 它给出规则，
而你需要遵循这些规则才能继续。

另一个回答也许更令人满意。
这个领域的人们一直在努力发展
威力与一般性之间的恰当平衡。
这个定义的设计，使得它包含了证明线性组合空间
一切有趣而重要性质所需的条件。
随着我们的推进，我们将从定义给出的条件出发，
导出线性组合的集合所自然具有的全部性质。

下面的结果就是一例。
我们无需把这些性质纳入向量空间的定义，
因为它们可以由其中已列出的性质推出。

\begin{lemma} \label{lm:ElementaryPropertiesOfVectorSpaces}
%<*lm:ElementaryPropertiesOfVectorSpaces>
在任何向量空间 \( V \) 中，
对任意 \( \vec{v}\in V \) 与 \( r\in\Re \)，有
(1)~\( 0\cdot\vec{v}=\zero \)，
(2)~\( (-1\cdot\vec{v})+\vec{v}=\zero \)，以及
(3)~\( r\cdot\zero=\zero \)。
%</lm:ElementaryPropertiesOfVectorSpaces>
\end{lemma}

\begin{proof}
%<*pf:ElementaryPropertiesOfVectorSpaces0>
对于 (1)，注意
\( \vec{v}=(1+0)\cdot\vec{v}=\vec{v}+(0\cdot\vec{v}) \)。
在两边加上 \( \vec{v} \) 的加法逆元，
即满足 \( \vec{w}+\vec{v}=\zero \) 的向量 \( \vec{w} \)。
\begin{align*}
  \vec{w}+\vec{v}
  &=\vec{w}+\vec{v}+0\cdot\vec{v}  \\
  \zero
  &=\zero+0\cdot\vec{v}                   \\
  \zero
  &=0\cdot\vec{v}
\end{align*}
条目~(2) 很容易：
\(  (-1\cdot\vec{v})+\vec{v}=(-1+1)\cdot\vec{v}=0\cdot\vec{v}=\zero \)。
对于 (3)，
\( r\cdot\zero
  =
  r\cdot(0\cdot\zero)
  =
  (r\cdot 0)\cdot\zero
  =
  \zero \)
即可。
%</pf:ElementaryPropertiesOfVectorSpaces0>
\end{proof}

第二条
表明我们可以把 \( \vec{v} \) 的加法逆元
写作“\( -\vec{v}\, \)”，而不必担心
与 \( (-1)\cdot\vec{v} \) 混淆。

\medskip
回顾一下：
第一章中对高斯消去法的研究
使我们去考虑线性组合的集合。
因此在本章中我们把向量空间定义为
可以形成这种组合的结构，
% expressions of the form \( c_1\cdot\vec{v}_1+\dots+c_n\cdot\vec{v}_n \)
并对加法与标量
乘法运算施加一些简单条件。
一言以蔽之：~向量空间是
研究线性性的恰当语境。

% Finally, a comment.
由于它构成了完整的一章，尤其因为这还是第一章，
读者可能会猜想本书的目的
就是研究线性方程组。
事实是：我们并非在用向量空间来研究线性方程组，
而是让线性方程组
引领我们开始对向量空间的研究。
本小节中各式各样的例子表明，对向量空间的研究
本身就有趣且重要。
线性方程组不会消失。
但从今以后，我们主要的研究对象将是向量空间。

\begin{exercises}
  \item
    写出下列各向量空间中的零向量。
    \begin{exparts}
      \partsitem 在通常运算下，三次多项式构成的空间。
      \partsitem \( \nbym{2}{4} \) 矩阵构成的空间。
      \partsitem 空间
        \( \set{\map{f}{[0..1]}{\Re}\suchthat f\text{ 连续}} \)。
      \partsitem 一个自然数自变量的实值函数构成的空间。
    \end{exparts}
    \begin{answer}
      \begin{exparts}
        \partsitem \( 0+0x+0x^2+0x^3 \)
        \partsitem \( \begin{mat}[r]
                   0  &0  &0  &0  \\
                   0  &0  &0  &0
                 \end{mat} \)
        \partsitem 常值函数 \( f(x)=0 \)
        \partsitem 常值函数 \( f(n)=0 \)
      \end{exparts}
    \end{answer}
  \recommended \item
    求该向量在此向量空间中的加法逆元。
    \begin{exparts}
      \partsitem 在 \( \polyspace_3 \) 中，向量 \( -3-2x+x^2 \)。
      \partsitem 在空间 \( \nbyn{2} \) 中，
        \begin{equation*}
          \begin{mat}[r]
            1  &-1  \\
            0  &3
          \end{mat}.
        \end{equation*}
      \partsitem 在 \( \set{ae^x+be^{-x}\suchthat a,b\in\Re} \)——这是实变量
        \( x \) 的函数在通常运算下构成的空间——中，向量 \( 3e^x-2e^{-x} \)。
    \end{exparts}
    \begin{answer}
      \begin{exparts*}
        \partsitem \( 3+2x-x^2 \)
        \partsitem \( \begin{mat}[r]
                   -1  &+1  \\
                    0  &-3
                 \end{mat} \)
        \partsitem \( -3e^x+2e^{-x} \)
      \end{exparts*}
    \end{answer}
  \recommended \item
    对每一小题，列出三个元素，然后证明它是向量空间。
    \begin{exparts}
      \partsitem 一次多项式的集合
        \( \polyspace_1=\set{a_0+a_1x\suchthat a_0,a_1\in\Re} \)，
        运算为通常的多项式加法与标量乘法。
      \partsitem 一次多项式的集合
        \( \set{a_0+a_1x\suchthat a_0-2a_1=0} \)，
        运算为通常的多项式加法与标量乘法。
   \end{exparts}
   \textit{提示。}  以 \nearbyexample{PlaneThruOriginSubsp} 作为指导。
   十个条件中的大多数都只是验证而已。
   \begin{answer}
      \begin{exparts}
        \partsitem
          三个元素是：
          $1+2x$、$2-1x$ 与 $x$。
          （当然，可以有许多不同的答案。）

          验证过程与 \nearbyexample{ex:RealVecSpaces} 完全类似。
          我们先处理
          \nearbydefinition{def:VecSpace} 中与加法有关的
          条件 $1$-$5$。
          关于对加法封闭，即条件~(1)，
          注意当 $a+bx,c+dx\in\polyspace_1$ 时，
          $(a+bx)+(c+dx)=(a+c)+(b+d)x$ 是系数为实数的一次多项式，
          因而是
          $\polyspace_1$ 的元素。
          条件~(2) 的验证如下：当 $a+bx,c+dx\in\polyspace_1$ 时，
          $(a+bx)+(c+dx)=(a+c)+(b+d)x$，而按另一顺序相加则得
          $(c+dx)+(a+bx)=(c+a)+(d+b)x$，由于这些都是实数，
          故 $a+c=c+a$
          且 $b+d=d+b$。
          条件~(3) 类似：设
          $a+bx,c+dx,e+fx\in\polyspace$，则
          $((a+bx)+(c+dx))+(e+fx)=(a+c+e)+(b+d+f)x$
          而
          $(a+bx)+((c+dx)+(e+fx))=(a+c+e)+(b+d+f)x$，两者相等
          （也就是说，实数加法满足结合律，故 $(a+c)+e=a+(c+e)$
          且 $(b+d)+f=b+(d+f)$）。
          对于条件~(4)，注意一次多项式
          $0+0x\in\polyspace_1$ 满足
          $(a+bx)+(0+0x)=a+bx$ 与
          $(0+0x)+(a+bx)=a+bx$。
          对于这一部分的最后一个条件，即条件~(5)，
          注意对任意 $a+bx\in\polyspace_1$，其加法逆元
          是 $-a-bx\in\polyspace_1$，因为
          $(a+bx)+(-a-bx)=(-a-bx)+(a+bx)=0+0x$。

          接下来我们还要检查与
          标量乘法有关的条件~(6)-(10)。
          对于~(6)，即空间对标量乘法封闭这一条件，
          设 $r$ 是实数，$a+bx$ 是
          $\polyspace_1$ 的元素，则 $r(a+bx)=(ra)+(rb)x$ 是
          $\polyspace_1$ 的元素，因为它是系数为
          实数的一次多项式。
          条件~(7) 成立，因为
          $(r+s)(a+bx)=r(a+bx)+s(a+bx)$ 可由实数乘法的分配律
          得到。
          条件~(8) 类似：
          $r((a+bx)+(c+dx))=r((a+c)+(b+d)x)=r(a+c)+r(b+d)x=(ra+rc)+(rb+rd)x
          =r(a+bx)+r(c+dx)$。
          对于~(9)，我们有
          $(rs)(a+bx)=(rsa)+(rsb)x=r(sa+sbx)=r(s(a+bx))$。
          最后，条件~(10) 是
          $1(a+bx)=(1a)+(1b)x=a+bx$。
        \partsitem
          将该集合记为 $P$。
          本练习的前一小题对系数没有限制，
          而这里我们只关注满足 $a_0-2a_1=0$ 的一次多项式，
          即常数项减去一次项系数的两倍为零的一次多项式。
          于是，$P$ 的三个典型元素是
          $2+1x$、$6+3x$ 与 $-4-2x$。

          对于条件~(1)，我们必须证明：
          若两个满足该限制的一次多项式相加，
          则所得仍是一次满足该限制的一次多项式：这里的论证是，
          若 $a+bx,c+dx\in P$，则
          $(a+bx)+(c+dx)=(a+c)+(b+d)x$ 是 $P$ 的元素，因为
          $(a+c)-2(b+d)=(a-2b)+(c-2d)=0+0=0$。
          条件~(2) 可以这样验证：
          当 $a+bx,c+dx\in\polyspace_1$ 时，
          $(a+bx)+(c+dx)=(a+c)+(b+d)x$，而按另一顺序相加则得
          $(c+dx)+(a+bx)=(c+a)+(d+b)x$，由于这些都是实数，
          故 $a+c=c+a$
          且 $b+d=d+b$。
          （也就是说，这一条件不受该限制的影响，
          其验证与本练习第一小题中的验证相同。）
          条件~(3) 同样不受这一附加限制的影响：
          设
          $a+bx,c+dx,e+fx\in\polyspace$，则
          $((a+bx)+(c+dx))+(e+fx)=(a+c+e)+(b+d+f)x$
          而
          $(a+bx)+((c+dx)+(e+fx))=(a+c+e)+(b+d+f)x$，两者相等。
          对于条件~(4)，注意一次多项式
          $0+0x\in P$ 满足该限制，因为
          其常数项减去其一次项系数的两倍
          为零，于是本小题中的
          验证即可适用：
          $0+0x\in\polyspace_1$ 满足
          $(a+bx)+(0+0x)=a+bx$ 与
          $(0+0x)+(a+bx)=a+bx$。
          为检查条件~(5)，
          注意对任意 $a+bx\in P$，其加法逆元
          是 $-a-bx$，因为
          它是 $P$ 的元素（由 $a+bx\in P$ 知
          $a-2b=0$，两边同乘 $-1$ 得 $-a+2b=0$），
          并且如第一小题那样，它起到加法逆元的作用：
          $(a+bx)+(-a-bx)=(-a-bx)+(a+bx)=0+0x$。

          我们还必须检查
          标量乘法的条件~(6)-(10)。
          对于~(6)，即空间对标量乘法封闭这一条件，
          设 $r$ 是实数且 $a+bx\in P$（于是 $a-2b=0$），
          则 $r(a+bx)=(ra)+(rb)x$ 是
          $P$ 的元素，因为它是系数为实数、
          且满足 $(ra)-2(rb)=r(a-2b)=0$ 的一次多项式。
          条件~(7) 成立的理由与本练习第一小题中
          相同，因为
          $(r+s)(a+bx)=r(a+bx)+s(a+bx)$ 可由实数乘法的分配律
          得到。
          条件~(8) 也与第一小题中相同：
          $r((a+bx)+(c+dx))=r((a+c)+(b+d)x)=r(a+c)+r(b+d)x=(ra+rc)+(rb+rd)x
          =r(a+bx)+r(c+dx)$。
          对~(9) 也是如此：
          $(rs)(a+bx)=(rsa)+(rsb)x=r(sa+sbx)=r(s(a+bx))$。
          最后，条件~(10) 也是如此：
          $1(a+bx)=(1a)+(1b)x=a+bx$。
       \end{exparts}
     \end{answer}
  \item
    对每一小题，列出三个元素，然后证明它是向量空间。
    \begin{exparts}
      \partsitem 实元素的 \( \nbyn{2} \) 矩阵的集合，运算为通常的矩阵运算。
      \partsitem 实元素且 $2,1$ 元为零的 \( \nbyn{2} \) 矩阵的集合，
        运算为通常的矩阵运算。
    \end{exparts}
    \begin{answer}
      可参照
      \nearbyexample{ex:RealVecSpaces} 作为指导。
      （\textit{评注。}
      由于许多条件都很容易检查，
      有时人们会觉得一定是漏掉了什么。
      请记住：容易做或常规，与不必做是两回事。）
      \begin{exparts}
        \partsitem
          以下是三个元素。
          \begin{equation*}
            \begin{mat}
              1  &2  \\
              3  &4
            \end{mat},\,
            \begin{mat}
              -1  &-2  \\
              -3  &-4
            \end{mat},\,
            \begin{mat}
              0  &0  \\
              0  &0
            \end{mat}
          \end{equation*}

          对于~(1)，$\nbyn{2}$ 实矩阵的和是 $\nbyn{2}$
          实矩阵。
          对于~(2)，我们考察两个矩阵的和
          \begin{equation*}
            \begin{mat}
              a  &b  \\
              c  &d
            \end{mat}+
            \begin{mat}
              e  &f  \\
              g  &h
            \end{mat}
            =
            \begin{mat}
              a+e  &b+f  \\
              c+g  &d+h
            \end{mat}
          \end{equation*}
          并利用实数加法的交换律
          \begin{equation*}
            =\begin{mat}
              e+a  &f+b  \\
              g+c  &h+d
            \end{mat}
            =
            \begin{mat}
              e  &f  \\
              g  &h
            \end{mat}
            +
            \begin{mat}
              a  &b  \\
              c  &d
            \end{mat}
          \end{equation*}
          来验证矩阵加法的交换性。
          条件~(3) 的验证，即矩阵加法的结合律，
          与前面的验证类似：
          \begin{equation*}
            \big(
              \begin{mat}
                a  &b  \\
                c  &d
              \end{mat}
              +\begin{mat}
                e  &f  \\
                g  &h
              \end{mat}
            \big)
            +
            \begin{mat}
              i  &j  \\
              k  &l
            \end{mat}
            =
            \begin{mat}
              (a+e)+i  &(b+f)+j  \\
              (c+g)+k  &(d+h)+l
            \end{mat}
          \end{equation*}
          而
          \begin{equation*}
            \begin{mat}
              a  &b  \\
              c  &d
            \end{mat}
            +
              \big(
              \begin{mat}
                e  &f  \\
                g  &h
              \end{mat}
              +
              \begin{mat}
                i  &j  \\
                k  &l
              \end{mat}
            \big)
            =
            \begin{mat}
              a+(e+i)  &b+(f+j)  \\
              c+(g+k)  &d+(h+l)
            \end{mat}
          \end{equation*}
          且由于实数加法满足结合律，两者逐项相同。
          对于~(4)，这个空间的零元素是全为零的 $\nbyn{2}$
          矩阵。
          条件~(5) 成立，因为对任意 $\nbyn{2}$ 矩阵~$A$，
          其加法逆元是各元素为 $A$ 相应元素的相反数的矩阵，
          即矩阵 $-1\cdot A$。

          条件~($6$)
          成立，因为 $\nbyn{2}$ 矩阵的标量倍
          仍是 $\nbyn{2}$ 矩阵。
          对于条件~(7)，我们有：
          \begin{multline*}
            (r+s)
            \begin{mat}
              a  &b  \\
              c  &d
            \end{mat}
            =
            \begin{mat}
              (r+s)a  &(r+s)b  \\
              (r+s)c  &(r+s)d
            \end{mat}               \\
            =
            \begin{mat}
              ra+sa  &rb+sb  \\
              rc+sc  &rd+sd
            \end{mat}
            =
            r
            \begin{mat}
              a  &b  \\
              c  &d
            \end{mat}
            +
            s
            \begin{mat}
              a  &b  \\
              c  &d
            \end{mat}
          \end{multline*}
          条件~(8) 的处理方式相同。
          \begin{multline*}
            r
            \big(
              \begin{mat}
                a  &b  \\
                c  &d
              \end{mat}
              +
              \begin{mat}
                e  &f  \\
                g  &h
              \end{mat}
            \big)
            =
            r
            \begin{mat}
              a+e  &b+f  \\
              c+g  &d+h
            \end{mat}
            =
            \begin{mat}
              ra+re  &rb+rf  \\
              rc+rg  &rd+rh
            \end{mat}              \\
            =
            r
            \begin{mat}
              a  &b  \\
              c  &d
            \end{mat}
            +
            r
            \begin{mat}
              e  &f  \\
              g  &h
            \end{mat}
            =
            r
            \big(
            \begin{mat}
              a  &b  \\
              c  &d
            \end{mat}
            +
            \begin{mat}
              e  &f  \\
              g  &h
            \end{mat}
            \big)
          \end{multline*}
          对于~(9)，我们有：
          \begin{equation*}
            (rs)
            \begin{mat}
              a  &b  \\
              c  &d
            \end{mat}
            =
            \begin{mat}
              rsa  &rsb  \\
              rsc  &rsd
            \end{mat}
            =
            r
            \begin{mat}
              sa  &sb  \\
              sc  &sd
            \end{mat}
            =
            r\big( s
            \begin{mat}
              a  &b  \\
              c  &d
            \end{mat}
            \big)
          \end{equation*}
          条件~(10) 同样简单。
          \begin{equation*}
            1
            \begin{mat}
              a  &b  \\
              c  &d
            \end{mat}
            =
            \begin{mat}
              1\cdot a  &1\cdot b  \\
              1\cdot c  &1\cdot d
            \end{mat}
            =
            \begin{mat}
              sa  &sb  \\
              sc  &sd
            \end{mat}
          \end{equation*}
        \partsitem
          本小题与本练习前一小题的差别仅在于：
          我们限定于第二行第一列
          为零的矩阵集合 $T$。
          以下是 $T$ 的三个元素。
          \begin{equation*}
            \begin{mat}
              1  &2  \\
              0  &4
            \end{mat},\,
            \begin{mat}
              -1  &-2  \\
              0  &-4
            \end{mat},\,
            \begin{mat}
              0  &0  \\
              0  &0
            \end{mat}
          \end{equation*}
          本小题的某些验证与本练习第一小题相同，
          下面我们只做那些不同的验证。

          对于~(1)，$2,1$ 元为零的 $\nbyn{2}$ 实矩阵之和
          也是 $2,1$ 元为零的 $\nbyn{2}$ 实矩阵。
          \begin{equation*}
            \begin{mat}
              a  &b  \\
              0  &d
            \end{mat}
            +
            \begin{mat}
              e  &f  \\
              0  &h
            \end{mat}
            \begin{mat}
              a+e  &b+f  \\
              0    &d+h
            \end{mat}
          \end{equation*}
          前一小题中给出的条件~(2) 的验证
          在本小题同样适用。
          条件~(3) 也是如此。
          对于~(4)，注意全为零的 $\nbyn{2}$
          矩阵是 $T$ 的元素。
          条件~(5) 成立，因为对任意 $\nbyn{2}$ 矩阵~$A$，
          其加法逆元是矩阵
          $-1\cdot A$，因此 $2,1$ 元为零的矩阵的加法逆元
          也是 $2,1$ 元为零的矩阵。

          条件~$6$ 成立，因为 $2,1$ 元为零的 $\nbyn{2}$ 矩阵的标量倍
          仍是 $2,1$ 元为零的 $\nbyn{2}$ 矩阵。
          条件~(7) 的验证与前一小题相同。
          条件~(8)、(9) 与~(10) 的验证也是如此。
      \end{exparts}
    \end{answer}
  \recommended \item
    对每一小题，列出三个元素，然后证明它是向量空间。
    \begin{exparts}
      \partsitem 三个分量的行向量的集合，运算为其通常运算。
      \partsitem 集合
        \begin{equation*}
          \set{\colvec{x \\ y \\ z \\ w}\in\Re^4\suchthat x+y-z+w=0}
        \end{equation*}
        运算取从 $\Re^4$ 继承的运算。
    \end{exparts}
    \begin{answer}
      % Most of the conditions are easy to check; use
      % \nearbyexample{ex:RealVecSpaces} as a guide.
      \begin{exparts}
        \partsitem
          三个元素是 $\rowvec{1  &2  &3}$、
          $\rowvec{2  &1  &3}$ 与 $\rowvec{0  &0  &0}$。

          我们必须检查 \nearbydefinition{def:VecSpace} 中的条件 (1)-(10)。
          条件 (1)-(5) 与加法有关。
          对于条件~(1)，回顾两个三分量行向量的和
          \begin{equation*}
            \rowvec{a  &b  &c}
            +\rowvec{d  &e  &f}
            =\rowvec{a+d  &b+e  &c+f}
          \end{equation*}
          也是三分量行向量
         （字母 $a,\ldots,f$ 全都表示实数）。
         对~(2) 的验证是常规的
          \begin{multline*}
            \rowvec{a  &b  &c}
            +\rowvec{d  &e  &f}
            =\rowvec{a+d  &b+e  &c+f}   \\
            =\rowvec{d+a  &e+b  &f+c}
            =\rowvec{d  &e  &f}
            +\rowvec{a  &b  &c}
          \end{multline*}
          （第二个等号成立是因为三个分量都是实数，
          而实数加法可交换）。
         条件~(3) 的验证类似。
          \begin{multline*}
            \big(\rowvec{a  &b  &c}
              +\rowvec{d  &e  &f}
            \big)
            +\rowvec{g  &h  &i}
            =\rowvec{(a+d)+g  &(b+e)+h  &(c+f)+i}    \\
            =\rowvec{a+(d+g)  &b+(e+h)  &c+(f+i)}
            =\rowvec{a  &b  &c}
            +\big(\rowvec{d  &e  &f}
               +\rowvec{g  &h  &i}
             \big)
          \end{multline*}
          对于~(4)，注意三分量行向量
          $\rowvec{0  &0  &0}$ 是加法单位元：
          $\rowvec{a  &b  &c}+\rowvec{0  &0  &0}=\rowvec{a  &b  &c}$。
          为验证条件~(5)，设给定集合中的
          元素 $\rowvec{a  &b  &c}$，并注意
          $\rowvec{-a  &-b  &-c}$ 也在此集合中且具有
          所需性质：
          $\rowvec{a  &b  &c}+\rowvec{-a  &-b  &-c}=\rowvec{0  &0  &0}$。

          条件 (6)-(10) 涉及标量乘法。
          为验证~(6)，即空间对所给的
          标量乘法运算封闭，
          注意 $r\rowvec{a  &b  &c}=\rowvec{ra  &rb  &rc}$ 是
          实分量的三分量行向量。
          对于~(7)，我们计算如下。
          \begin{multline*}
            (r+s)\rowvec{a  &b  &c}=\rowvec{(r+s)a  &(r+s)b  &(r+s)c}
            =\rowvec{ra+sa  &rb+sb  &rc+sc}                             \\
            =\rowvec{ra  &rb  &rc}+\rowvec{sa  &sb  &sc}
            =r\rowvec{a  &b  &c}+s\rowvec{a  &b  &c}
          \end{multline*}
          条件~(8) 非常类似。
          \begin{multline*}
            r\big(\rowvec{a  &b  &c}+\rowvec{d  &e  &f}\big)  \\
            \begin{aligned}
            &=r\rowvec{a+d  &b+e  &c+f}
            =\rowvec{r(a+d)  &r(b+e)  &r(c+f)}                     \\
            &=\rowvec{ra+rd  &rb+re  &rc+rf}
            =\rowvec{ra  &rb  &rc}+\rowvec{rd  &re  &rf}          \\
            &=r\rowvec{a  &b  &c}+r\rowvec{d  &e  &f}
            \end{aligned}
          \end{multline*}
          条件~(9) 的计算也是如此。
          \begin{equation*}
            (rs)\rowvec{a  &b  &c}
             =\rowvec{rsa  &rsb  &rsc}
             =r\rowvec{sa  &sb  &sc}
             =r\big(s\rowvec{a  &b  &c}\big)
          \end{equation*}
          条件~(10) 同样是常规的：
          $1\rowvec{a  &b  &c}
             =\rowvec{1\cdot a  &1\cdot b  &1\cdot c}
             =\rowvec{a  &b  &c}$。
        \partsitem
          将该集合记为 $L$。
          加法的封闭性，即条件~(1)，需要检查：若各加数都是~$L$ 的成员，则
          和
          \begin{equation*}
            \colvec{a \\ b \\ c \\ d}
            +\colvec{e \\ f \\ g \\ h}
            =
            \colvec{a+e \\ b+f \\ c+g \\ d+h}
          \end{equation*}
          也是 \( L \) 的成员；这是成立的，因为它满足
          $L$ 的成员资格判据：
          \( (a+e)+(b+f)-(c+g)+(d+h)
          =(a+b-c+d)+(e+f-g+h)=0+0 \)。
          条件~(2)、(3) 与~(5) 的验证与本练习第一部分中的
          验证类似。
          对于条件~(4)，注意零向量是~$L$ 的成员，
          因为其第一分量加第二分量、减第三分量、
          再加第四分量，总和为零。

          条件~(6)，即标量乘法的封闭性，类似：
          当该向量是~$L$ 的元素时，
          \begin{equation*}
            r\colvec{a  \\ b \\ c \\ d}
            =\colvec{ra  \\  rb  \\  rc  \\  rd}
          \end{equation*}
          也是~$L$ 的元素，因为 $ra+rb-rc+rd=r(a+b-c+d)=r\cdot 0=0$。
          条件~(7)、(8)、(9) 与~(10) 的验证与本练习
          前一小题中相同。
     \end{exparts}
     \end{answer}
  \recommended \item \label{exer:NotVectorSpaces}
    证明下列每一个都不是向量空间。
    （\textit{提示。}  为检查封闭性，列出每个集合的两个成员，
    并对它们尝试一些运算。）
    \begin{exparts}
      \partsitem 在从 \( \Re^3 \) 继承的运算下，集合
        \begin{equation*}
          \set{\colvec{x \\ y \\ z}\in\Re^3\suchthat x+y+z=1}
        \end{equation*}
      \partsitem 在从 \( \Re^3 \) 继承的运算下，集合
        \begin{equation*}
          \set{\colvec{x \\ y \\ z}\in\Re^3\suchthat x^2+y^2+z^2=1}
        \end{equation*}
      \partsitem 在通常的矩阵运算下，
        \begin{equation*}
          \set{\begin{mat}
                 a  &1  \\
                 b  &c
               \end{mat} \suchthat a,b,c\in\Re}
        \end{equation*}
      \partsitem 在通常的多项式运算下，
        \begin{equation*}
          \set{a_0+a_1x+a_2x^2\suchthat a_0,a_1,a_2\in\Re^+}
        \end{equation*}
        其中 $\Re^+$ 是大于零的实数构成的集合
      \partsitem 在继承的运算下，
        \begin{equation*}
          \set{\colvec{x \\ y}\in\Re^2\suchthat
               \text{\( x+3y=4 \) 且 \( 2x-y=3 \) 且 \( 6x+4y=10 \)} }
        \end{equation*}
    \end{exparts}
    \begin{answer}
      每一小题中都将该集合记为 \( Q \)。
      对其中某些小题，还有其他正确方法可证明 $Q$ 不是
      向量空间。
      \begin{exparts}
        \partsitem 它对加法不封闭；它不满足
          条件~(1)。
          \begin{equation*}
            \colvec[r]{1 \\ 0 \\ 0},
            \colvec[r]{0 \\ 1 \\ 0}\in Q
            \qquad
            \colvec[r]{1 \\ 1 \\ 0}\not\in Q
          \end{equation*}
        \partsitem 它对加法不封闭。
          \begin{equation*}
            \colvec[r]{1 \\ 0 \\ 0},
            \colvec[r]{0 \\ 1 \\ 0}\in Q
            \qquad
            \colvec[r]{1 \\ 1 \\ 0}\not\in Q
          \end{equation*}
        \partsitem 它对加法不封闭。
          \begin{equation*}
            \begin{mat}[r]
              0  &1  \\
              0  &0
            \end{mat},
            \,
            \begin{mat}[r]
              1  &1  \\
              0  &0
            \end{mat}\in Q
            \qquad
            \begin{mat}[r]
              1  &2  \\
              0  &0
            \end{mat}\not\in Q
          \end{equation*}
        \partsitem 它对标量乘法不封闭。
          \begin{equation*}
            1+1x+1x^2\in Q
            \qquad
            -1\cdot(1+1x+1x^2)\not\in Q
          \end{equation*}
        \item 它是空集，违反条件~(4)。
      \end{exparts}
    \end{answer}
  \item
    定义加法与标量乘法运算，
    使复数成为 \( \Re \) 上的向量空间。
    \begin{answer}
      通常的运算
      \( (v_0+v_1i)+(w_0+w_1i)=(v_0+w_0)+(v_1+w_1)i \) 与
      \( r(v_0+v_1i)=(rv_0)+(rv_1)i \) 即可满足要求。
      检查很容易。
    \end{answer}
  \item
    在通常的加法与标量乘法
    运算下，有理数集是 \( \Re \) 上的向量空间吗？
    \begin{answer}
       不是，它对标量乘法不封闭，因为例如
       \( \pi\cdot (1) \) 不是有理数。
    \end{answer}
  \item
    证明
    变量 \( x,y,z \) 的线性组合构成的集合
    在自然的加法与标量乘法
    运算下是向量空间。
    \begin{answer}
      自然运算是
      \( (v_1x+v_2y+v_3z)+(w_1x+w_2y+w_3z)=(v_1+w_1)x+(v_2+w_2)y+(v_3+w_3)z \)
      与 \( r\cdot(v_1x+v_2y+v_3z)=(rv_1)x+(rv_2)y+(rv_3)z \)。
      检查它是向量空间很容易；以
      \nearbyexample{ex:RealVecSpaces} 作为指导。
    \end{answer}
  \item
    证明
    这不是向量空间：实元素的
    二行列向量的集合，其运算如下。
    \begin{equation*}
      \colvec{x_1 \\ y_1}
      +\colvec{x_2 \\ y_2}
      =\colvec{x_1-x_2 \\ y_1-y_2}
      \qquad
      r\cdot\colvec{x \\ y}
      =\colvec{rx \\ ry}
    \end{equation*}
    \begin{answer}
      “\( + \)” 运算不可交换（即条件~(2) 不
      满足）；很容易找出该集合的两个成员来
      证实这一断言。
    \end{answer}
  \item
    证明或反驳：\( \Re^3 \) 在下列
    运算下是向量空间。
    \begin{exparts}
      \partsitem \(
               \colvec{x_1 \\ y_1 \\ z_1}
               +\colvec{x_2 \\ y_2 \\ z_2}
               =\colvec{0 \\ 0 \\ 0}
               \quad\text{且}\quad
               r\colvec{x \\ y \\ z}
               =\colvec{rx \\ ry \\ rz} \)
      \partsitem \(
               \colvec{x_1 \\ y_1 \\ z_1}
               +\colvec{x_2 \\ y_2 \\ z_2}
               =\colvec{0 \\ 0 \\ 0}
               \quad\text{且}\quad
               r\colvec{x \\ y \\ z}
               =\colvec{0 \\ 0 \\ 0} \)
    \end{exparts}
    \begin{answer}
      \begin{exparts}
        \partsitem 它不是向量空间。
          \begin{equation*}
            (1+1)\cdot\colvec[r]{1 \\ 0 \\ 0}\neq
            \colvec[r]{1 \\ 0 \\ 0}
            +\colvec[r]{1 \\ 0 \\ 0}
          \end{equation*}
        \partsitem 它不是向量空间。
          \begin{equation*}
            1\cdot\colvec[r]{1 \\ 0 \\ 0}\neq\colvec[r]{1 \\ 0 \\ 0}
          \end{equation*}
      \end{exparts}
    \end{answer}
  \recommended \item
    对每一小题，判断它是否为向量空间；
    所指的运算是自然的那些。
    \begin{exparts}
      \partsitem \definend{对角} \( \nbyn{2} \) 矩阵
        \begin{equation*}
          \set{\begin{mat}
            a  &0  \\
            0  &b
          \end{mat}\suchthat a,b\in\Re}
        \end{equation*}
      \partsitem 这个 \( \nbyn{2} \) 矩阵集合
        \begin{equation*}
          \set{\begin{mat}
            x    &x+y  \\
            x+y  &y
          \end{mat}\suchthat x,y\in\Re}
        \end{equation*}
      \partsitem 这个集合
        \begin{equation*}
          \set{\colvec{x \\ y \\ z \\ w}\in\Re^4
               \suchthat x+y+z+w=1}
        \end{equation*}
      \partsitem 函数的集合
        \( \set{\map{f}{\Re}{\Re}\suchthat df/dx+2f=0} \)
      \partsitem 函数的集合
        \( \set{\map{f}{\Re}{\Re}\suchthat df/dx+2f=1} \)
    \end{exparts}
    \begin{answer}
      对每个“是”的答案，只需在该空间中检查几个
      线性组合即可。
      对每个“否”的答案，给出其中一个
      条件不成立的具体例子。
      \begin{exparts}
        \partsitem 是。
          随手取的一个线性组合如下。
          \begin{equation*}
            2\cdot\begin{mat}
              3  &0 \\
              0  &4
            \end{mat}
            +5\cdot\begin{mat}
              6  &0 \\
              0  &7
            \end{mat}
          \end{equation*}
          这个集合成为向量空间的关键在于：
          结果是对角矩阵。
          \begin{equation*}
            =\begin{mat}
              36 &0 \\
              0  &43
            \end{mat}
          \end{equation*}
        \partsitem 是。
          线性组合如下。
          \begin{equation*}
            10\cdot\begin{mat}
              1  &3  \\
              3  &2
            \end{mat}
            +11\cdot\begin{mat}
              3  &7  \\
              7  &4
            \end{mat}
          \end{equation*}
          我们看到结果保持了如下形式：左上
          与右下元素相加得到右上元素，
          同时也得到左下元素。
          \begin{equation*}
            =\begin{mat}
               43  &107  \\
              107  &64
            \end{mat}
          \end{equation*}
        \partsitem 否，这个集合对自然加法
          运算不封闭。
          全部分量都是 $1/4$ 的向量是这个集合的成员，
          但它自加所得的结果，即
          全部分量都是 $1/2$ 的向量，不是成员。
        \partsitem 是。
        \partsitem 否，\( f(x)=e^{-2x}+(1/2) \) 在集合中，
           但 \( 2\cdot f \) 不在（即条件~(6) 不成立）。
      \end{exparts}
    \end{answer}
  \recommended \item
    证明或反驳：这是一个向量空间——满足 \( f(7)=0 \) 的
    单实变量实值函数 \( f \) 构成的集合。
    \begin{answer}
      它是向量空间。
      向量空间定义中的大多数条件的验证都是常规的；这里我们
      只检查封闭性。
      对于加法，
      \( (f_1+f_2)\,(7)=f_1(7)+f_2(7)=0+0=0 \)。
      对于标量乘法，
      \( (r\cdot f)\,(7)=rf(7)=r0=0 \)。
    \end{answer}
  \recommended \item
    证明：正实数集 \( \Re^+ \)
    是向量空间，其中我们把 “\( x+y \)” 解释为
    \( x \) 与 \( y \) 的乘积（于是 \( 2+3 \) 为 \( 6 \)），
    把 “\( r\cdot x \)” 解释为 \( x \) 的 \( r \) 次幂。
    \begin{answer}
      我们逐条检查 \nearbydefinition{def:VecSpace}。

      首先，对`\( + \)'封闭
      成立，因为两个正实数之积是
      正实数。
      第二个条件满足，因为实数乘法可交换。
      类似地，由于实数乘法可结合，第三条也成立。
      对于第四个条件，注意用
      \( 1\in\Re^+ \)中的数去乘一个数不会改变它。
      第五，任何正实数都有一个倒数，它是正实数。

      第六条，对`\( \cdot \)'封闭， 
      成立是因为正实数的任意次幂仍是
      正实数。
      第七个条件就是\( v^{r+s} \)等于
      \( v^r \)与\( v^s \)之积这条法则。
      第八个条件说\( (vw)^r=v^rw^r \)。
      第九个条件断言\( (v^r)^s=v^{rs} \)。
      最后一个条件说\( v^1=v \)。  
    \end{answer}
  \item 
      \( \set{(x,y)\suchthat x,y\in\Re} \)在这些
      运算下是向量空间吗？
      \begin{exparts}
        \partsitem \( (x_1,y_1)+(x_2,y_2)=(x_1+x_2,y_1+y_2) \)
        以及\( r\cdot (x,y)=(rx,y) \)
        \partsitem \( (x_1,y_1)+(x_2,y_2)=(x_1+x_2,y_1+y_2) \)
        以及\( r\cdot (x,y)=(rx,0) \)
      \end{exparts}
      \begin{answer}
        \begin{exparts}
           \partsitem 不是：\( 1\cdot(0,1)+1\cdot(0,1)\neq (1+1)\cdot(0,1) \)。
           \partsitem 不是；与前一答案相同的计算表明，向量空间定义中的
              一个条件被违反了。 
              违反向空间条件的另一个例子是
              \( 1\cdot (0,1)\neq (0,1) \)。 
        \end{exparts}  
      \end{answer}
  \item 
    证明或否证：这是
    一个向量空间——次数大于或等于二的多项式
    连同零多项式构成的集合。
    \begin{answer}
      它不是向量空间，因为它对加法不封闭，例如
      \( (x^2)+(1+x-x^2) \)不在此集合中。  
    \end{answer}
  \item 
    此时“相同”还只是一种
    直观，但尽管如此，请对每个向量空间指出使得
    该空间与\( \Re^k \)“相同”的\( k \)。
    \begin{exparts}
      \partsitem 通常运算下的\( \nbym{2}{3} \)矩阵
      \partsitem \( \nbym{n}{m} \)矩阵（在它们通常的运算下）
      \partsitem \( \nbyn{2} \)矩阵的这个集合
        \begin{equation*}
          \set{\begin{mat}
            a  &0  \\
            b  &c
          \end{mat} \suchthat a,b,c\in\Re}
        \end{equation*}
      \partsitem \( \nbyn{2} \)矩阵的这个集合
        \begin{equation*}
          \set{\begin{mat}
            a  &0  \\
            b  &c
          \end{mat} \suchthat a+b+c=0}
        \end{equation*}
    \end{exparts}
    \begin{answer}
      \begin{exparts}
        \partsitem \( 6 \)
        \partsitem \( nm \)
        \partsitem \( 3 \)
        \partsitem 为看出答案是\( 2 \)，把它改写为
        \begin{equation*}
          \set{\begin{mat}
            a  &0  \\
            b  &-a-b
          \end{mat} \suchthat a,b\in\Re}
        \end{equation*}
        可见其中有两个参数。
      \end{exparts}  
     \end{answer}
  \item 
    用\( \vec{+} \)表示向量加法、
    用\( \,\vec{\cdot}\, \)表示标量乘法，
    重新表述向量空间的定义。
    \begin{answer}
      {
      \def\plus{\mathbin{\vec{+}}}
      \def\tim{\mathbin{\vec{\cdot}}}
        \definend{向量空间}\index{vector space!definition}
        （\( \Re \)上的）由一个集合\( V \)连同两个运算`\( \plus \)'与
        `\( \tim \)'组成，满足以下条件。
        当\( \vec{v},\vec{w}\in V \)时， 
        (1)~它们的\definend{向量和}
          \( \vec{v}\plus\vec{w} \)是\( V \)中的一个元素。
        若\( \vec{u},\vec{v},\vec{w}\in V \)，则 
        (2)~\( \vec{v}\plus\vec{w}=\vec{w}\plus\vec{v} \)且 
        (3)~\( (\vec{v}\plus\vec{w})\plus\vec{u}
                 =\vec{v}\plus(\vec{w}\plus\vec{u}) \)。
        (4)~存在一个\definend{零向量}
          \( \zero\in V \)，使得
          对所有\( \vec{v}\in V\)都有\( \vec{v}\plus\zero=\vec{v}\, \)。
        (5)~每个\( \vec{v}\in V \)都有一个
          \definend{加法逆元}
          \( \vec{w}\in V \)，使得\( \vec{w}\plus\vec{v}=\zero \)。
        若\( r,s \)是\definend{标量\/}
        （即\( \Re \)的成员），
        且\( \vec{v},\vec{w}\in V \)，则 
        (6)~每个
           \definend{标量倍数}
           \( r\cdot\vec{v} \)都在\( V \)中。
        若\( r,s\in\Re \)且\( \vec{v},\vec{w}\in V \)，则
        (7)~\( (r+ s)\cdot\vec{v}=r\cdot\vec{v}\plus s\cdot\vec{v} \)， 
        以及(8)~\( r\tim (\vec{v}+\vec{w})
           =r\tim\vec{v}+r\tim\vec{w} \)，
        以及(9)~\( (rs)\tim \vec{v} =r\tim (s\tim\vec{v}) \)，
        以及(10)~\( 1\tim \vec{v}=\vec{v} \)。
     }
    \end{answer}
  \item 
    证明这些结论。
    \begin{exparts}
      \partsitem 对任意$\vec{v}\in V$，若$\vec{w}\in V$是$\vec{v}$的加法 
        逆元，则$\vec{v}$是
        $\vec{w}$的加法逆元。
        因此，一个向量是它自己的任意一个加法逆元的加法逆元。
      \partsitem 向量加法可从左边消去：~若 
        \( \vec{v},\vec{s},\vec{t}\in V \)，
        则\( \vec{v}+\vec{s}=\vec{v}+\vec{t}\, \)蕴涵
        \( \vec{s}=\vec{t} \)。
    \end{exparts}
    \begin{answer}
      \begin{exparts}
        \partsitem 设\( V \)是向量空间， 
          设\( \vec{v}\in V \)，并
          假设\( \vec{w}\in V \)是$\vec{v}$的加法逆元，
          于是\( \vec{w}+\vec{v}=\zero \)。
          因为加法可交换，
          \( \zero=\vec{w}+\vec{v}=\vec{v}+\vec{w} \)，
          所以\( \vec{v} \)也是 
          \( \vec{w} \)的加法逆元。
        \partsitem 设\( V \)是向量空间且设
          \( \vec{v},\vec{s},\vec{t}\in V \)。
          \( \vec{v} \)的加法逆元是\( -\vec{v} \)，于是
          \( \vec{v}+\vec{s}=\vec{v}+\vec{t} \)给出
          \( -\vec{v}+\vec{v}+\vec{s}=-\vec{v}+\vec{v}+\vec{t} \)，
          即\( \zero+\vec{s}=\zero+\vec{t} \)，从而
          \( \vec{s}=\vec{t} \)。
      \end{exparts}  
     \end{answer}
  \item 
    向量空间的定义并未明确说出
    \( \zero+\vec{v}=\vec{v} \) 
    （它说的是\( \vec{v}+\zero=\vec{v} \)）。
    证明它在任何向量空间中仍必成立。
    \begin{answer}
      加法可交换，所以在任何向量空间中，
      对任意向量\( \vec{v} \)我们有
      \( \vec{v}=\vec{v}+\zero=\zero+\vec{v} \)。  
    \end{answer}
  \recommended \item
    证明或否证：这是
    一个向量空间——通常运算下的全体矩阵之集。
    \begin{answer}
      它不是向量空间，因为尺寸不同的两个矩阵相加没有定义，
      因而该集合不满足封闭性
      条件。
    \end{answer}
  \item 
    向量空间中每个元素都有一个加法逆元。
    某些元素能否有两个或更多个？
    \begin{answer}
      向量空间中每个元素有且只有一个加法
      逆元。

      设\( V \)是向量空间且设\( \vec{v}\in V \)。
      若\( \vec{w}_1,\vec{w}_2\in V \)都是
      \( \vec{v} \)的加法逆元，则考虑\( \vec{w}_1+\vec{v}+\vec{w}_2 \)。
      一方面，我们有它等于$\vec{w}_1+(\vec{v}+\vec{w}_2)=
      \vec{w}_1+\zero=\vec{w}_1$。
      另一方面，我们有它等于$(\vec{w}_1+\vec{v})+\vec{w}_2=
      \zero+\vec{w}_2=\vec{w}_2$。
      因此$\vec{w}_1=\vec{w}_2$。
    \end{answer}
  \item 
    \begin{exparts}
      \partsitem 证明\( \Re^3 \)中每个过原点的点、直线或平面 
         在继承的运算下是向量空间。
      \partsitem 若它不包含原点呢？
    \end{exparts}
    \begin{answer}
     \begin{exparts}
      \partsitem 每个这样的集合都具有形式
        \( \set{r\cdot\vec{v}+s\cdot\vec{w}\suchthat r,s\in\Re} \)，
        其中\( \vec{v},\vec{w} \)之一或两者可以是\( \zero \)。
        在继承的运算下，加法的封闭性
        \( (r_1\vec{v}+s_1\vec{w})+(r_2\vec{v}+s_2\vec{w})
           =(r_1+r_2)\vec{v}+(s_1+s_2)\vec{w} \)
        与标量乘法的封闭性
        \( c(r\vec{v}+s\vec{w})=(cr)\vec{v}+(cs)\vec{w} \)
        是容易的。
        其他条件也是常规的。
      \partsitem 这样的集合在继承的
        运算下不能是向量空间，因为它没有零元素。
     \end{exparts}  
    \end{answer}
  \item 
    利用向量空间的想法，我们可以轻松地重新证明
    齐次线性方程组的解集要么
    有一个元素，要么有无穷多个元素。 
    设\( \vec{v}\in V \)不是\( \zero \)。
    \begin{exparts}
      \partsitem 证明\( r\cdot\vec{v}=\zero \)当且仅当\( r=0 \)。
      \partsitem 证明\( r_1\cdot\vec{v}=r_2\cdot\vec{v} \)
      当且仅当\( r_1=r_2 \)。
      \partsitem 证明任何非平凡的向量空间都是无限的。
      \partsitem 利用齐次线性方程组的非空解集是向量空间这一事实
        来得出结论。
    \end{exparts}
    \begin{answer}
      设\( \vec{v}\in V \)不是\( \zero \)。
      \begin{exparts}
        \partsitem 当且仅当的一个方向是清楚的：~若$r=0$，
          则$r\cdot\vec{v}=\zero$。
          另一个方向，设\( r \)是非零标量。
          若\( r\vec{v}=\zero \)，则
          \( (1/r)\cdot r\vec{v}=(1/r)\cdot \zero \)表明
          $\vec{v}=\zero$，与假设矛盾。
        \partsitem 当\( r_1,r_2 \)是标量时， 
          \( r_1\vec{v}=r_2\vec{v}\, \)
          成立当且仅当\( (r_1-r_2)\vec{v}=\zero \)。
          由前一项，于是\( r_1-r_2=0 \)。
        \partsitem 非平凡空间有向量 
          \( \vec{v}\neq\zero \)。
          考虑集合\( \set{k\cdot\vec{v}\suchthat k\in\Re} \)。
          由前一项，这个集合是无限的。
        \partsitem 解集要么是平凡的，要么是非平凡的。
          在第二种情形，它是无限的。   
     \end{exparts}  
    \end{answer}
  \item 
    在自然运算下这是向量空间吗：一个实变量的
    可微实值函数之集？
    \begin{answer}
      是。
      第一学期微积分的一个定理说，可微函数之和可微，且
      \( (f+g)^\prime=f^\prime+g^\prime \)；还说 
      可微函数的倍数
      可微，且\( (r\cdot f)^\prime=r\,f^\prime \)。 
    \end{answer}
  \item 
    \definend{复数上的向量空间}\index{vector space!complex scalars}%
    \index{complex numbers!vector space over}
    $\C$的定义
    与实数上的向量空间相同，只是标量取自
    \( \C \)而非\( \Re \)。
    证明下列每一个都是复数上的向量空间。
    （回忆复数如何相加与相乘：
    \( (a_0+a_1i)+(b_0+b_1i)=(a_0+b_0)+(a_1+b_1)i \)以及
    \( (a_0+a_1i)(b_0+b_1i)=(a_0b_0-a_1b_1)+(a_0b_1+a_1b_0)i \)。）
    \begin{exparts}
      \partsitem 系数为复数的二次多项式之集
      \partsitem 这个集合
        \begin{equation*}
          \set{\begin{mat}
                 0  &a  \\
                 b  &0
               \end{mat}\suchthat a,b\in\C\text{\ 且\ }
                               a+b=0+0i }
        \end{equation*}
    \end{exparts}
    \begin{answer}
      这个验证是常规的。
      注意`\( 1 \)'是\( 1+0i \)，零元素是下面这些。
      \begin{exparts}
        \partsitem \( (0+0i)+(0+0i)x+(0+0i)x^2 \)
        \partsitem \( \begin{mat}
                   0+0i  &0+0i  \\
                   0+0i  &0+0i
                 \end{mat} \)
      \end{exparts}  
    \end{answer}
  \item 
    指出所有\( \Re^n \)共有、
    但未被列为
    向量空间要求的一条性质。
    \begin{answer}
      向量空间的定义中显著缺失的
      是距离的度量。  
    \end{answer}
  \item 
    \begin{exparts}
      \partsitem 证明对任意四个向量
        \( \vec{v}_1,\ldots,\vec{v}_4\in V \)，我们可以
        以任何方式结合它们的和而不改变结果。
        \begin{multline*}
          ((\vec{v}_1+\vec{v}_2)+\vec{v}_3)+\vec{v}_4
          =(\vec{v}_1+(\vec{v}_2+\vec{v}_3))+\vec{v}_4  
          =(\vec{v}_1+\vec{v}_2)+(\vec{v}_3+\vec{v}_4)  \\
          =\vec{v}_1+((\vec{v}_2+\vec{v}_3)+\vec{v}_4)  
          =\vec{v}_1+(\vec{v}_2+(\vec{v}_3+\vec{v}_4))
        \end{multline*}
        这使我们可以写成
        `\( \vec{v}_1+\vec{v}_2+\vec{v}_3+\vec{v}_4 \)'而无歧义。
      \partsitem 证明对任意多个向量的和，任何两种结合方式
        给出相同的和。
        （\textit{提示。}  对向量的个数使用归纳法。）
    \end{exparts}
    \begin{answer}
      \begin{exparts}
        \partsitem 小小的重新排列即可奏效。
          \begin{align*}
            (\vec{v}_1+(\vec{v}_2+\vec{v}_3))+\vec{v}_4
            &=((\vec{v}_1+\vec{v}_2)+\vec{v}_3)+\vec{v}_4  \\
            &=(\vec{v}_1+\vec{v}_2)+(\vec{v}_3+\vec{v}_4)  \\
            &=\vec{v}_1+(\vec{v}_2+(\vec{v}_3+\vec{v}_4))  \\
            &=\vec{v}_1+((\vec{v}_2+\vec{v}_3)+\vec{v}_4)
          \end{align*}
          上面每个等号都来自三个向量的结合性，
          它作为一条条件写在向量空间的定义中。
          例如，第二个`$=$'应用法则
          $(\vec{w}_1+\vec{w}_2)+\vec{w}_3=\vec{w}_1+(\vec{w}_2+\vec{w}_3)$，
          取$\vec{w}_1$为$\vec{v}_1+\vec{v}_2$， 
          取$\vec{w}_2$为$\vec{v}_3$，
          取$\vec{w}_3$为$\vec{v}_4$。 
        \partsitem 归纳的基础是三个向量的情形。
          这一情形
          \( \vec{v}_1+(\vec{v}_2+\vec{v}_3)
          =(\vec{v}_1+\vec{v}_2)+\vec{v}_3 \)
          是向量空间定义中的条件之一。

          归纳步骤：假设三个向量的任意两个和、
          四个向量的任意两个和、
          \ldots、$k$个向量的任意两个和， 
          无论我们如何给这些和加括号都相等。
          我们将证明任意\( k+1 \)个向量的和都等于这一个
          \( ((\cdots((\vec{v}_1+\vec{v}_2)+\vec{v}_3)+\cdots)+\vec{v}_k)
              +\vec{v}_{k+1} \)。

          任何加了括号的和都有一个最外层的`\( + \)'。
          设它位于\( \vec{v}_m \)与\( \vec{v}_{m+1} \)之间，
          于是这个和看起来像这样。
          \begin{equation*}
            (\cdots\,\vec{v}_1\cdots\vec{v}_m\,\cdots)
           +(\cdots\,\vec{v}_{m+1}\cdots\vec{v}_{k+1}\,\cdots)
          \end{equation*}
          后一半涉及的加法少于$k+1$次，所以
          由归纳假设我们可以重新加括号，
          使得它从里到外自左向右读，特别地，
          使得其最外层的`$+$'恰好出现在 
          $\vec{v}_{k+1}$之前。
          \begin{equation*}
            =(\cdots\,\vec{v}_1\,\cdots\,\vec{v}_m\,\cdots)
             +((\cdots(\vec{v}_{m+1}+\vec{v}_{m+2})+\cdots+\vec{v}_{k})
                 +\vec{v}_{k+1})
          \end{equation*}
          应用三个对象之和的结合性
          \begin{equation*}
            =((\,\cdots\, \vec{v}_1\,\cdots\,\vec{v}_m\,\cdots\,)
             +(\,\cdots\,(\vec{v}_{m+1}+\vec{v}_{m+2})+\cdots\,\vec{v}_k))
             +\vec{v}_{k+1}
          \end{equation*}
          最后在这些最外层括号的内部应用归纳假设，即完成证明。
      \end{exparts}  
    \end{answer}
  \item \nearbyexample{ex:ColsIntEntNotVS}给出$\Re^2$的一个子集，
   在显然的运算下它不是向量空间，因为
   虽然它对加法封闭，却对标量
   乘法不封闭。
   考虑平面上分量同号或为~$0$的向量组成的集合。
   证明这个集合对标量乘法封闭但对 
   加法不封闭。
   \begin{answer}
     设$\vec{v}$是$\Re^2$的成员，分量为$v_1$与$v_2$。
     两个分量同号或为~$0$这一条件可以缩写
     为$v_1v_2\geq 0$。

     为证这个集合对标量乘法封闭，注意
     $r\vec{v}$的分量满足$(rv_1)(rv_2)=r^2(v_1v_2)$，
     且$r^2\geq 0$，所以$r^2v_1v_2\geq 0$。

     为证这个集合对加法不封闭，我们只需给出一个
     例子。  
     分量为$-1$与~$0$的向量加到
     分量为$0$与~$1$的向量上，得到一个分量为$-1$与~$1$、
     符号混合的向量。
   \end{answer}
  % \item 
  %   For any vector space, a subset that is itself a vector space
  %   under the inherited operations
  %   (e.g., a plane through the origin inside of \( \Re^3 \))
  %   is a \definend{subspace}.
  %   \begin{exparts}
  %     \partsitem Show that \( \set{a_0+a_1x+a_2x^2\suchthat a_0+a_1+a_2=0} \)
  %       is a subspace of the vector space of degree~two polynomials.
  %     \partsitem Show that this is a subspace of the \( \nbyn{2} \) matrices.
  %       \begin{equation*}
  %         \set{\begin{mat}
  %                a  &b  \\
  %                c  &0
  %              \end{mat} \suchthat a+b=0}
  %       \end{equation*}
  %     \partsitem Show that a nonempty subset \( S \) of a real vector
  %       space is a
  %       subspace if and only if it is closed under linear combinations of
  %       pairs of vectors: whenever \( c_1,c_2\in\Re \) and
  %       \( \vec{s}_1,\vec{s}_2\in S \) then the combination
  %       \( c_1\vec{v}_1+c_2\vec{v}_2 \) is in \( S \).
  %   \end{exparts}
  %   \begin{answer}
  %     \begin{exparts}
  %       \partsitem We outline the check of the conditions from
  %         \nearbydefinition{def:VecSpace}.

  %         Additive closure holds because if \( a_0+a_1+a_2=0 \)
  %         and \( b_0+b_1+b_2=0 \) then
  %         \begin{equation*}
  %           (a_0+a_1x+a_2x^2)+(b_0+b_1x+b_2x^2)=
  %           (a_0+b_0)+(a_1+b_1)x+(a_2+b_2)x^2
  %         \end{equation*}
  %         is in the set since
  %         \( (a_0+b_0)+(a_1+b_1)+(a_2+b_2)=(a_0+a_1+a_2)+(b_0+b_1+b_2) \)
  %         is zero.
  %         The second through fifth conditions are easy.

  %         Closure under scalar multiplication holds because
  %         if \( a_0+a_1+a_2=0 \) then
  %         \begin{equation*}
  %           r\cdot(a_0+a_1x+a_2x^2)=
  %           (ra_0)+(ra_1)x+(ra_2)x^2
  %         \end{equation*}
  %         is in the set as \( ra_0+ra_1+ra_2=r(a_0+a_1+a_2) \) is zero.
  %         The remaining conditions here are also easy.
  %       \partsitem This is similar to the prior answer.
  %       \partsitem Call the vector space \( V \).
  %         We have two implications: left to right, if \( S \) is a subspace
  %         then it is closed under linear combinations of pairs of vectors and,
  %         right to left, if a nonempty subset is closed under linear
  %         combinations of pairs of vectors then it is a subspace.
  %         The left to right implication is easy; we here sketch the
  %         other one by
  %         assuming \( S \) is nonempty and closed, and checking the conditions
  %         of \nearbydefinition{def:VecSpace}.

  %         First, to show closure under addition, if
  %         \( \vec{s}_1,\vec{s}_2\in S \) then \( \vec{s}_1+\vec{s}_2\in S \)
  %         as \( \vec{s}_1+\vec{s}_2=1\cdot\vec{s}_1+1\cdot\vec{s}_2 \).
  %         Second, for any \( \vec{s}_1,\vec{s}_2\in S \), because addition
  %         is inherited from \( V \), the sum \( \vec{s}_1+\vec{s}_2 \)
  %         in \( S \) equals the sum \( \vec{s}_1+\vec{s}_2 \)
  %         in \( V \) and that equals the sum \( \vec{s}_2+\vec{s}_1 \) in
  %         \( V \) and that in turn equals the sum \( \vec{s}_2+\vec{s}_1 \) in
  %         \( S \).
  %         The argument for the third condition is similar to that for the
  %         second.
  %         For the fourth, suppose that 
  %         \( \vec{s} \) is in the nonempty set \( S \)
  %         and note that \( 0\cdot\vec{s}=\zero\in S \); showing that 
  %         the \( \zero \) of
  %         \( V \) acts under the inherited operations as the additive identity
  %         of \( S \) is easy.
  %         The fifth condition is satisfied because for any \( \vec{s}\in S \)
  %         closure under linear combinations shows that the vector
  %         \( 0\cdot\zero+(-1)\cdot\vec{s} \) is in \( S \); showing 
  %         that it is the
  %         additive inverse of \( \vec{s} \) under the inherited operations is
  %         routine.

  %         The proofs for the remaining conditions are similar.
  %     \end{exparts}  
  %   \end{answer}
\end{exercises}





















 
\subsection{子空间与张成集}
\index{subspace|(}
在\nearbyexample{PlaneThruOriginSubsp}中，我们见到了一个向量空间，它是 $\Re^2$ 的一个子集：一条过原点的直线。
在那里，向量空间 $\Re^2$ 的内部还包含着另一个
向量空间，即这条直线。

\begin{definition} \label{df:Subspace}
%<*df:Subspace>
对于任意向量空间，
\definend{子空间}\index{vector space!subspace}\index{subspace!definition}
是指在继承的运算之下，其本身也是一个向量空间的子集。
%</df:Subspace>
\end{definition}

\begin{example}  \label{ex:PlaneSubspRThree}
这个过原点的平面
\begin{equation*}
  P=\set{\colvec{x \\ y \\ z}\suchthat x+y+z=0}
\end{equation*}
是 \( \Re^3 \) 的一个子空间。
正如定义所要求的，
该平面的运算继承自更大的空间，
也就是说，
向量在 $P$ 中的加法与在 $\Re^3$ 中的加法一样
\begin{equation*}
   \colvec{x_1 \\ y_1 \\ z_1}+\colvec{x_2 \\ y_2 \\ z_2}
   =\colvec{x_1+x_2 \\ y_1+y_2 \\ z_1+z_2}
\end{equation*}
且标量乘法也与 $\Re^3$ 中的相同。
为证明 $P$ 是子空间，我们只需注意它是一个子集，
然后验证它是一个空间。
我们不必检验全部十个条件，只检验两个封闭性条件。
关于对加法封闭，注意若
两个加项分别满足
$x_1+y_1+z_1=0$ 与 $x_2+y_2+z_2=0$，则其和满足
$(x_1+x_2)+(y_1+y_2)+(z_1+z_2)=(x_1+y_1+z_1)+(x_2+y_2+z_2)=0$。
关于对标量乘法封闭，若
$x+y+z=0$，则其标量倍满足
$rx+ry+rz=r(x+y+z)=0$。
\end{example}

\begin{example}   \label{ex:SubspacesRTwo}
\( \Re^2 \) 中的 \( x \) 轴
是一个子空间，其中
加法与标量乘法运算都是
继承而来的那些运算。
\begin{equation*}
  \colvec{x_1 \\ 0}
    +
  \colvec{x_2 \\ 0}
    =
  \colvec{x_1+x_2 \\ 0}
  \qquad
  r\cdot\colvec{x \\ 0}
  =\colvec{rx \\ 0}
\end{equation*}
与前面的例题一样，要从定义直接验证它是子空间，
我们只需注意它是一个
子集，然后检验它满足
向量空间定义中的各个条件。
例如两个封闭性条件都
满足：第二个分量为零的两个向量相加，所得向量的第二个分量仍为零；
标量乘以第二个分量为零的向量，
所得向量的第二个分量仍为零。
\end{example}

\begin{example}
$\Re^2$ 的另一个子空间是
它的平凡子空间。
\begin{equation*}
  \set{\colvec[r]{0 \\ 0}}
\end{equation*}
\end{example}

%<*ProperSubspace>
任何向量空间都有一个平凡子空间
\( \set{\zero\,} \)。\index{subspace!trivial}%
\index{trivial subspace}
与之相反的极端是，任何向量空间都以自身作为自己的一个子空间。
不是整个空间的子空间称为
\definend{真}\index{subspace!proper}%
\index{proper subspace} 子空间。
%</ProperSubspace>

\begin{example}  \label{ex:LinSubspPolyThree}
不是 $\Re^n$ 的向量空间也有子空间。
三次多项式的空间
\( \set{a+bx+cx^2+dx^3\suchthat a,b,c,d\in\Re} \)
有一个由全部一次多项式
\( \set{m+nx\suchthat m,n\in\Re} \)
组成的子空间。
\end{example}

\begin{example}
另一个不是 $\Re^n$ 的子集的子空间的例子出现在向量空间的定义之后。
\nearbyexample{ex:RealValuedFcns}中全体一元实值函数构成的空间
\( \set{f\suchthat \map{f}{\Re}{\Re} } \) 含有\nearbyexample{ex:SpaceSatisfyingDiffQ}中满足
限制 $(d^2\,f/dx^2)+f=0$ 的那些函数组成的子空间。
\end{example}

\begin{example} \label{ex:OperNotInherit}
该定义要求
加法与标量乘法运算
必须是继承自更大空间的那两种运算。
集合 \( S=\set{1} \) 是 \( \Re^1 \) 的一个子集。
而且在运算 $1+1=1$ 与 $r\cdot 1=1$ 之下，
集合 $S$ 是一个向量空间，确切地说，是一个平凡空间。
然而 $S$ 不是 \( \Re^1 \) 的子空间，因为那些运算不是
继承的运算，毕竟 \( \Re^1 \) 中当然有 \( 1+1=2 \)。
\end{example}

\begin{example}  \label{cex:RPlusNotSubSp}
子空间本身是向量空间，因此必须满足封闭性
条件。
集合 \( \Re^+ \) 不是向量空间 \( \Re^1 \) 的子空间，
因为在继承的运算下它对标量乘法不封闭：
若 \( \vec{v}=1 \)，则 \( -1\cdot\vec{v}\not\in\Re^+ \)。
\end{example}

下面的结果表明，\nearbyexample{cex:RPlusNotSubSp} 是典型的。
一个子集若非空且使用继承的运算，
则它不成其为子空间的唯一可能就是不封闭。

\begin{lemma}     \label{th:SubspIffClosed} \index{subspace!closed}
%<*lm:SubspIffClosed>
对于向量空间的一个非空子集 \( S \)，在继承的
运算之下，下列诸命题
等价。\appendrefs{equivalence of statements}   %\spacefactor=1000
\begin{tfae}
  \item \( S \) 是该向量空间的一个子空间
  \item \( S \) 对两个向量的线性组合封闭：
    对任意向量 \( \vec{s}_1,\vec{s}_2\in S \) 与标量 \( r_1,r_2 \)，
    向量 \( r_1\vec{s}_1+r_2\vec{s}_2 \) 属于 \( S \)
  \item \( S \) 对任意多个向量的线性组合封闭：
    对任意向量 \( \vec{s}_1,\ldots,\vec{s}_n\in S \) 与标量
    \( r_1, \ldots,r_n \)，
    向量 \( r_1\vec{s}_1+\cdots+r_n\vec{s}_n \) 是 \( S \) 的元素。
\end{tfae}
%</lm:SubspIffClosed>
\end{lemma}
\noindent 简言之，一个子集是
子空间当且仅当它对线性组合封闭。

\begin{proof}
%<*pf:SubspIffClosed0>
“下列诸命题等价”是指每一对
命题都等价。
\begin{equation*}
  (1)\!\iff\!(2)
  \qquad
  (2)\!\iff\!(3)
  \qquad
  (3)\!\iff\!(1)
\end{equation*}
我们将通过确立
\( (1)\implies (3)\implies (2)\implies (1) \) 来证明这一等价性。
采取这一策略的理由是如下观察：蕴含关系
\( (1)\implies (3) \) 与 \( (3)\implies (2) \) 都很容易，于是我们只需
论证 \( (2)\implies (1) \)。
%</pf:SubspIffClosed0>

%<*pf:SubspIffClosed1>
设 \( S \) 是向量空间
$V$ 的一个非空子集，且对两个向量的组合封闭。
我们将通过检验各项条件来证明 $S$ 是一个向量空间。

向量空间的定义中有五条关于加法的条件。
第一，关于对加法封闭，若
\( \vec{s}_1,\vec{s}_2\in S \)，则 \( \vec{s}_1+\vec{s}_2\in S \)，
因为它是
两个向量的组合，
而我们正假设~\( S \) 对这种组合封闭。
第二，对任意 \( \vec{s}_1,\vec{s}_2\in S \)，由于加法
继承自 \( V \)，
\( S \) 中的和 \( \vec{s}_1+\vec{s}_2 \)
等于 \( V \) 中的和 \( \vec{s}_1+\vec{s}_2 \)，
而后者等于
\( V \) 中的和 \( \vec{s}_2+\vec{s}_1 \)（因为 $V$ 是向量空间，其加法可交换），
进而后者又等于 \( S \) 中的和 \( \vec{s}_2+\vec{s}_1 \)。
第三个条件的论证与第二个类似。
第四，考虑 \( V \) 的零向量，注意
$S$ 对两个向量的线性组合的封闭性给出：
\( 0\cdot\vec{s}+0\cdot\vec{s}=\zero \) 是 \( S\) 的元素
（其中 \( \vec{s} \) 是非空集 \( S \) 的任一成员）；
容易检验 \( \zero \) 在继承的运算下充当 \( S \) 的加法
单位元。
第五个条件满足，因为对任意 \( \vec{s}\in S \)，
对两个向量的线性组合的封闭性表明
\( 0\cdot\zero+(-1)\cdot\vec{s} \) 是~\( S \) 的元素，而且它显然是
\( \vec{s} \) 在继承的运算下的
加法逆元。
%</pf:SubspIffClosed1>
关于标量乘法诸条件的验证类似；
见\nearbyexercise{exer:SubspIffClosed}。
\end{proof}

以后我们通常通过检验一个子集
满足命题~(2) 来验证它是子空间。

\begin{remark}
在本章开头，我们把向量空间作为线性组合在其中“有意义”的集合加以介绍。
% The vector space definition has ten conditions but eight of them, the
% conditions not about closure, simply ensure that referring to the
% operations by the names `addition' and `scalar multiplication' is sensible.
% For instance, calling an operation `$+$' would not
% be odd unless it is commutative, as in condition~(2). 
% The proof above checks that the subspace satisfies these 
% eight non-closure conditions because they hold in the
% surrounding vector space, provided that the nonempty set $S$ satisfies
% \nearbylemma{th:SubspIffClosed}'s statement~(2) 
% (e.g., commutativity of addition in $S$ follows right from
% commutativity of addition in $V$).
% There is a second way in which we can regard vector spaces as structures in 
% which linear combinations are ``sensible.''
% It has to do with the closure conditions.
\nearbylemma{th:SubspIffClosed} 的诸命题 (1)-(3)
表明我们总能对形如
$r_1\vec{s}_1+r_2\vec{s}_2$
的表达式给出意义，
也就是说，所描述的向量在集合~$S$ 中。

作为对比，考虑由分量之和不小于零的二元向量组成的集合 $T$。
在这里我们不能随手写出任意的线性组合如 $2\vec{t}_1-3\vec{t}_2$
并确信其结果是 $T$ 的元素。
\end{remark}

\nearbylemma{th:SubspIffClosed} 提示我们，理解向量空间的一种好方式
是把它看作无限制线性组合的集合。
接下来的两个例子取一些空间，把对它们的描述改写成这种形式。

\begin{example}
要证明 $\Re^3$ 的这个过原点的平面子集
\begin{equation*}
  S=\set{\colvec{x \\ y \\ z}\suchthat x-2y+z=0}
\end{equation*}
在列向量通常的加法与标量乘法
运算下是子空间，我们可以检验它非空且对
两个向量的线性组合封闭。
但还有另一种方法。
把 $x-2y+z=0$ 看作只含一个方程的线性方程组并加以参数化：
将首变量用自由变量表示为 $x=2y-z$。
\begin{equation*}
     S
     =\set{\colvec{2y-z \\ y \\ z}\suchthat y,z\in\Re}
     =\set{y\colvec[r]{2 \\ 1 \\ 0}+
            z\colvec[r]{-1 \\ 0 \\ 1}\suchthat y,z\in\Re}
    \tag{$*$}
\end{equation*}
现在，为证明它是子空间，考虑
$r_1\vec{s}_1+r_2\vec{s}_2$。
每个 $\vec{s}_i$ 都是 ($*$) 中那两个向量的线性组合，
所以这是一个线性组合的线性组合。
\begin{equation*}
  r_1\cdot(y_1\colvec[r]{2 \\ 1 \\ 0}+
            z_1\colvec[r]{-1 \\ 0 \\ 1})
  +
  r_2\cdot(y_2\colvec[r]{2 \\ 1 \\ 0}+
            z_2\colvec[r]{-1 \\ 0 \\ 1})
\end{equation*}
线性组合引理（Lemma~One.III.\ref{lm:LinearCombinationLemma}）
表明总和是那两个向量的线性组合，
于是 \nearbylemma{th:SubspIffClosed} 的命题~(2) 得到满足。
\end{example}

\begin{example} \label{ex:ParamSubspace}
这是全体 \( \nbyn{2} \) 矩阵的空间 $\matspace_{\nbyn{2}}$ 的一个子空间。
\begin{equation*}
  L=\set{\begin{mat}
         a  &0  \\
         b  &c
       \end{mat}
       \suchthat a+b+c=0}
\end{equation*}
为参数化，把条件写成 $a=-b-c$。
\begin{equation*}
  L
  =\set{\begin{mat}
         -b-c  &0  \\
         b     &c
       \end{mat}
       \suchthat b,c\in\Re}
  =\set{b\begin{mat}[r]
         -1    &0  \\
         1     &0
       \end{mat}
       +c\begin{mat}[r]
         -1    &0  \\
         0     &1
       \end{mat}
       \suchthat b,c\in\Re}
\end{equation*}
如上所述，我们已把这个子空间描述为无限制线性
组合的集合。
为证明它是子空间，注意来自
$L$ 的向量的线性组合是线性组合的线性组合，因此命题~(2)
成立。
\end{example}

% Parametrization is easy, but important.
% We shall use it often.

\begin{definition} \label{df:Span}
%<*df:Span>
向量空间的一个非空子集 \( S \) 的\definend{张成}\index{span}\index{sets!span of}\index{closure}（或称
\definend{线性闭包}）是由 \( S \) 中向量的一切线性组合组成的集合。
\begin{equation*}
  \spanof{S} =\{ c_1\vec{s}_1+\cdots+c_n\vec{s}_n
            \suchthat c_1,\ldots, c_n\in\Re
            \text{\ 且\ } \vec{s}_1,\ldots,\vec{s}_n\in S \}
\end{equation*}
向量空间的空子集的张成是它的平凡子空间。
%</df:Span>
\end{definition}

%<*NotationForSpan>
\noindent 关于张成的记号没有完全标准的。
这里用的方括号很常见，
“$\mbox{span}(S)$” 与 “$\mbox{sp}(S)$” 也很常用。
%</NotationForSpan>

\begin{remark}
在第一章，当我们证明了可以把齐次线性方程组的解
集写成
$\set{c_1\vec{\beta}_1+\cdots+c_k\vec{\beta}_k\suchthat
  c_1,\ldots,c_k\in\Re}$
之后，我们把它描述为由诸 $\smash{\vec{\beta}}$ “生成”的集合。
现在我们称它为
$\set{\vec{\beta}_1,\ldots,\vec{\beta}_k}$ 的张成。

也应从那个证明中回忆起：
空集的张成定义为集合 \( \set{\zero} \)，
这是出于这样一个约定：平凡的线性组合，即零个向量的组合，其和为~\( \zero \)。
此外，把空集的张成定义为平凡子空间
也很方便，因为它使得像下一个结果那样的定理
无需为空集另立例外。
\end{remark}

\begin{lemma}   \label{le:SpanIsASubsp}
%<*lm:SpanIsASubsp>
在向量空间中，任何子集的张成都是子空间。
%</lm:SpanIsASubsp>
\end{lemma}

\begin{proof}
%<*pf:SpanIsASubsp>
若子集 \( S \)
为空，则按定义其张成是平凡
子空间。
若 \( S\) 非空，则由\nearbylemma{th:SubspIffClosed}，我们
只需检验张成 \( \spanof{S} \) 对两个元素的线性组合封闭。
取该张成中的一对向量，
\( \vec{v}=c_1\vec{s}_1+\cdots+c_n\vec{s}_n \) 与
\( \vec{w}=c_{n+1}\vec{s}_{n+1}+\cdots+c_m\vec{s}_m \)，
则线性组合
\begin{multline*}
  p\cdot(c_1\vec{s}_1+\cdots+c_n\vec{s}_n)+
       r\cdot(c_{n+1}\vec{s}_{n+1}+\cdots+c_m\vec{s}_m)  \\
  =
  pc_1\vec{s}_1+\cdots+pc_n\vec{s}_n
    +rc_{n+1}\vec{s}_{n+1}+\cdots+rc_m\vec{s}_m
\end{multline*}
是~\( S \) 中元素的线性组合，
因此是 \( \spanof{S} \) 的元素
（来自 $\vec{v}$ 的某些 $\vec{s}_i$ 可能与来自 $\vec{w}$ 的某些 $\vec{s}_j$ 相同，但这无关紧要）。
%</pf:SpanIsASubsp>
\end{proof}

这个引理的逆
也成立：任何子空间都是某个集合的张成，
因为子空间显然是它自身的张成，即由其全体成员
构成的集合的张成。
于是，向量空间的一个子集是子空间当且仅当它是一个张成。
这符合下述
直觉：理解向量空间的好方式是把它看作
线性组合在其中有意义的集合。

合起来，\nearbylemma{th:SubspIffClosed} 与
\nearbylemma{le:SpanIsASubsp} 表明：向量空间的子集 $S$ 的
张成是包含 $S$ 的全体成员的最小子空间。

\begin{example}   \label{ex:SpanSingVec}
在任意向量空间 \( V \) 中，对任意向量 \( \vec{v}\in V \)，集合
\( \set{r\cdot\vec{v} \suchthat r\in\Re} \) 是 \( V \) 的子空间。
例如，对任意向量 \( \vec{v}\in\Re^3 \)，
包含该向量的过原点直线
\( \set{k\vec{v}\suchthat k\in\Re } \) 是 \( \Re^3 \) 的子空间。
即使 $\vec{v}$ 是零向量，这也是对的，此时
它是退化直线，即平凡子空间。
\end{example}

\begin{example}
这个集合的张成
是整个 $\Re^2$。
\begin{equation*}
  \set{\colvec[r]{1 \\ 1},\colvec[r]{1 \\ -1}}
\end{equation*}
我们知道该张成是~$\Re^2$ 的某个子空间。
要检验它就是整个~$\Re^2$，
必须证明 $\Re^2$ 的任何成员都是这两个向量的
线性组合。
于是我们问：~对以实数 $x$ 与~$y$ 为分量的哪些向量，
存在标量 $c_1$ 与 $c_2$ 使下式成立？
\begin{equation*}
   c_1\colvec[r]{1 \\ 1}+c_2\colvec[r]{1 \\ -1}=\colvec{x \\ y}
  \tag{$*$}
\end{equation*}
高斯消元法
\begin{equation*}
  \begin{linsys}{2}
    c_1  &+  &c_2  &=  &x  \\
    c_1  &-  &c_2  &=  &y
  \end{linsys}
  \grstep{-\rho_1+\rho_2}
  \begin{linsys}{2}
    c_1  &+  &c_2    &=  &x\hfill  \\
         &   &-2c_2  &=  &-x+y
  \end{linsys}
\end{equation*}
经回代给出 $c_2=(x-y)/2$ 与 $c_1=(x+y)/2$。
这表明对任意 $x,y$ 都存在
适当的系数 $c_1,c_2$ 使~($*$) 成立\Dash 我们可以把 $\Re^2$ 的任意元素写成这两个
给定向量的线性组合。
例如，取 $x=1$ 与 $y=2$，系数 $c_2=-1/2$ 与
$c_1=3/2$ 即可。
\end{example}

既然张成是子空间，而我们又知道
理解子空间的好方式是
将其描述参数化，那么我们可以尝试用这种方式来理解一个集合的张成。

\begin{example}
在次数不超过二的多项式空间~\( \polyspace_2 \) 中，考虑
集合 \( S=\set{3x-x^2, 2x} \) 的张成。
由张成的定义，它是这两个向量的无限制线性组合
$\set{c_1(3x-x^2)+c_2(2x)\suchthat c_1,c_2\in\Re}$ 组成的集合。
显然这个张成中的多项式的常数项必为零。
这个必要条件也是充分的吗？

我们问的是：~对 $\polyspace_2$ 的哪些成员 $a_2x^2+a_1x+a_0$
存在 $c_1$ 与 $c_2$ 使
$a_2x^2+a_1x+a_0=c_1(3x-x^2)+c_2(2x)$？
系数相等的多项式相等，因此
我们要找 $a_2$、$a_1$ 与 $a_0$ 上的条件，
使这个三元组成为下面方程组的解。
\begin{equation*}
  \begin{linsys}{2}
    -c_1  &   &     &=  &a_2   \\
    3c_1  &+  &2c_2 &=  &a_1   \\
          &   &0    &=  &a_0
  \end{linsys}
\end{equation*}
高斯消元法与回代给出
$c_1=-a_2$，以及 $c_2=(3/2)a_2+(1/2)a_1$，以及 $0=a_0$。
于是
%the only condition on elements  $a_0+a_1x+a_2x^2$ of the span
%is the condition that we knew:
只要没有常数项，即 $a_0=0$，
我们就能给出系数 $c_1$ 与 $c_2$，
把该多项式描述为张成中的一个元素。
例如，对多项式 $0-4x+3x^2$，系数
$c_1=-3$ 与 $c_2=5/2$ 即可。
所以给定集合的张成是
$\spanof{S}=\set{a_1x+a_2x^2\suchthat a_1,a_2\in\Re}$。

顺带地，这表明
集合 \( \set{x,x^2} \) 张成同一个子空间。
一个空间可以有不止一个张成集。
张成这个子空间的另外两个集合是
\( \set{x,x^2,-x+2x^2} \) 与
\( \set{x,x+x^2,x+2x^2,\ldots\,} \)。
% (Usually we prefer to work with spanning sets that have only
% a small number of members.)
\end{example}

\begin{example}  \label{ex:SubspRThree}
下图展示了我们目前已知的 \( \Re^3 \) 的子空间：平凡
子空间、过原点的直线、
过原点的平面，以及整个空间。
（当然，图示只画出了无穷多种情形中的少数几个。
线段把子集与其
超集连接起来。）
在下一节我们将证明 $\Re^3$ 没有其他
类型的子空间，因此实际上这张图列出了全部子空间。

这把每个子空间
描述为成员个数最少的集合的张成。
据此，各子空间
自然分成一些层次\Dash 平面在同一层，
直线在另一层，
等等。
\begin{center}
  \setlength{\unitlength}{4pt}
  \begin{picture}(75,38)(0,-2) %subspaces of R-three
      \thinlines
      \put(45,31){\makebox(0,0)[bl]{
                        \scriptsize \( %\Re^3=
                                   \set{x\colvec[r]{1 \\ 0 \\ 0}
                                              +y\colvec[r]{0 \\ 1 \\ 0}
                                              +z\colvec[r]{0 \\ 0 \\ 1}} \)} }
    %Next the dimension two subspaces
      \put(43,31){\line(-4,-1){25} } % connects R3 to 2,a
      %set 2,a:
      \put(0,22){\makebox(0,0)[l]{\scriptsize\( \set{x\colvec[r]{1 \\ 0 \\ 0}
                                                 +y\colvec[r]{0 \\ 1 \\ 0} }\) }}
      \put(48,29){\line(-3,-1){12} } % connects R3 to 2,b
      %set 2,b:
      \put(20,22){\makebox(0,0)[l]{\scriptsize\( \set{x\colvec[r]{1 \\ 0 \\ 0}
                                                 +z\colvec[r]{0 \\ 0 \\ 1} }\) }}
      \put(53,29){\line(-1,-1){3}} % connects R3 to 2,c
      %set 2,c:
      \put(40,22){\makebox(0,0)[l]{\scriptsize\( \set{x\colvec[r]{1 \\ 1 \\ 0}
                                                 +z\colvec[r]{0 \\ 0 \\ 1} }\) }}
      \put(60,22){\makebox(0,0)[l]{$\cdots$} }
    %Next the dimension one subspaces
      \put(7,18){\line(-1,-4){1} } % connects 2,a to 1,a
      \put(22,18){\line(-3,-1){13} } % connects 2,c to 1,a
      %set 1,a:
      \put(0,10){\makebox(0,0)[l]{\scriptsize\( \set{x\colvec[r]{1 \\ 0 \\ 0}} \)} }
      \put(12,18){\line(1,-2){2} } % connects 2,a to 1,b
      %set 1,b:
      \put(12,10){\makebox(0,0)[l]{\scriptsize\( \set{y\colvec[r]{0 \\ 1 \\ 0}} \)} }
      \put(14,18){\line(2,-1){10} } % connects 2,a to 1,c
      %set 1,c:
      \put(24,10){\makebox(0,0)[l]{\scriptsize\( \set{y\colvec[r]{2 \\ 1 \\ 0}} \)} }
      \put(46,18){\line(-1,-1){3.25} } %connects 2,c to 1,d
      %set 1,d:
      \put(36,10){\makebox(0,0)[l]{\scriptsize\( \set{y\colvec[r]{1 \\ 1 \\ 1}} \)} }
      \put(48,10){\makebox(0,0)[l]{$\cdots$} }
    %Finally, the trivial subspace.
      \put(9,7.5){\line(4,-1){30} } %connects 1,a to trivial
      \put(18.9,7.7){\line(3,-1){20} } %connects 1,b to trivial
      \put(30.4,6.8){\line(2,-1){10} } %connects 1,c to trivial
      \put(41,6){\line(1,-2){1.5} } %connects 1,d to trivial
      \put(45,0){\makebox(0,0){\scriptsize\( \set{\colvec[r]{0 \\ 0 \\ 0} } \)} }
  \end{picture}
\end{center}
\end{example}

到本章目前为止我们已经看到：要研究
线性组合的性质，合适的框架是
对这些组合封闭的集合。
在第一小节中我们引入了这类集合，即向量空间，
并见到了丰富多样的例子。
在本小节中我们又见到了
更多的空间，它们是其他空间的子空间。
在所有这些多样性中有一个共同点。
上面的\nearbyexample{ex:SubspRThree}
把它揭示了出来：~向量空间与子空间最好被理解为张成，
尤其是少量向量的张成。
下一节将研究具有极小性的张成集。






\begin{exercises}
  \recommended \item
    下列 \( \nbyn{2} \) 矩阵构成的向量空间中的哪些子集，在继承的运算下是子空间？
    对每一个是子空间的，给出其参数化描述。
    对每一个不是的，给出一条不满足的条件。
    \begin{exparts}
      \partsitem \( \set{\begin{mat}
                      a  &0  \\
                      0  &b
                    \end{mat}  \suchthat a,b\in\Re}  \)
      \partsitem \( \set{\begin{mat}
                      a  &0  \\
                      0  &b
                    \end{mat}  \suchthat a+b=0} \)
      \partsitem \( \set{\begin{mat}
                      a  &0  \\
                      0  &b
                    \end{mat}  \suchthat a+b=5} \)
      \partsitem \( \set{\begin{mat}
                      a  &c  \\
                      0  &b
                    \end{mat}  \suchthat a+b=0, c\in\Re} \)
    \end{exparts}
    \begin{answer}
      根据 \nearbylemma{th:SubspIffClosed}，要判断 $\matspace_{\nbyn{2}}$
      的每个子集是否为子空间，只需检查它是否非空且封闭。
      \begin{exparts}
        \partsitem 是，容易验证它非空且封闭。
          这就是它的一个参数化描述。
          \begin{equation*}
            \set{a\begin{mat}[r]
                    1  &0  \\
                    0  &0
                  \end{mat}
                 +b\begin{mat}[r]
                    0  &0  \\
                    0  &1
                   \end{mat}
                 \suchthat a,b\in\Re}
          \end{equation*}
          顺便指出，这个参数化描述本身也表明它是一个子空间，
          因为它被表示为由两个矩阵构成的集合张成，
          而任何张成都是子空间。
        \partsitem 是；容易验证它非空且封闭。
         或者，如前一个答案所述，参数化描述的存在性
         表明它是一个子空间。
         要得到参数化描述，
         可将条件 $a+b=0$ 改写为 $a=-b$。
         于是有下式。
          \begin{equation*}
            \set{\begin{mat}
                    -b  &0  \\
                    0   &b
                  \end{mat}
                 \suchthat b\in\Re}
            =\set{b\begin{mat}[r]
                    -1  &0  \\
                    0   &1
                  \end{mat}
                 \suchthat b\in\Re}
          \end{equation*}
        \partsitem 不是。
          它对加法不封闭。
          例如，
          \begin{equation*}
            \begin{mat}[r]
              5  &0  \\
              0  &0
            \end{mat}
            +\begin{mat}[r]
              5  &0  \\
              0  &0
            \end{mat}
            =\begin{mat}[r]
              10  &0  \\
              0  &0
            \end{mat}
          \end{equation*}
          不在该集合中。
          （这个集合对标量乘法也不封闭，
          例如，它不包含零矩阵。）
        \partsitem 是。
          \begin{equation*}
            \set{b\begin{mat}[r]
                    -1  &0  \\
                    0   &1
                  \end{mat}
                 +c\begin{mat}[r]
                    0  &1  \\
                    0  &0
                   \end{mat}
                 \suchthat b,c\in\Re}
          \end{equation*}
      \end{exparts}
     \end{answer}
  \recommended \item
    这是 \( \polyspace_2 \) 的子空间吗：
    \( \set{a_0+a_1x+a_2x^2\suchthat a_0+2a_1+a_2=4} \)？
    如果是，给出其参数化描述。
    \begin{answer}
      不是，它不封闭。
      具体来说，它对标量乘法不封闭，因为它
      不包含零多项式。
    \end{answer}
  \item 下面的向量是否在该集合的张成之中？
    \begin{equation*}
      \colvec{1 \\ 0 \\ 3}
      \quad
      \set{\colvec{2 \\ 1 \\ -1},
           \colvec{1 \\ -1 \\ 1}}
    \end{equation*}
    \begin{answer}
      方程
      \begin{equation*}
        \colvec{1 \\ 0 \\ 3} =
                 c_1\colvec{2 \\ 1 \\ -1}
                 +c_2\colvec{1 \\ -1 \\ 1}
      \end{equation*}
      给出一个线性方程组
      \begin{equation*}
        \begin{amat}{2}
          2  &1  &1  \\
          1  &-1 &0  \\
          -1 &1  &3
        \end{amat}
        \grstep[(1/2)\rho_1+\rho_3]{(-1/2)\rho_1+\rho_2}
        \begin{amat}{2}
          2  &1    &1  \\
          0  &-3/2 &-1/2  \\
          0  &0  &3
        \end{amat}
      \end{equation*}
      该方程组无解，所以该向量不在张成之中。
    \end{answer}
  \recommended \item
    判断该向量是否张成于该集合的张成之中，所在
    空间如下。
    \begin{exparts}
      \partsitem \( \colvec[r]{2 \\ 0 \\ 1} \),
        \( \set{\colvec[r]{1 \\ 0 \\ 0},
                \colvec[r]{0 \\ 0 \\ 1}  } \),
        在 \( \Re^3 \) 中
      \partsitem \( x-x^3 \),
        \( \set{x^2,2x+x^2,x+x^3} \),
        在 \( \polyspace_3 \) 中
      \partsitem \( \begin{mat}[r]
                 0  &1  \\
                 4  &2
               \end{mat}  \),
        \( \set{\begin{mat}[r]
                  1  &0  \\
                  1  &1
                \end{mat},
                \begin{mat}[r]
                  2  &0  \\
                  2  &3
                \end{mat}  } \),
        在 \( \matspace_{\nbyn{2}} \) 中
    \end{exparts}
    \begin{answer}
      \begin{exparts}
         \partsitem 是，由
          \begin{equation*}
            r_1\colvec[r]{1 \\ 0 \\ 0}+r_2\colvec[r]{0 \\ 0 \\ 1}
              =\colvec[r]{2 \\ 0 \\ 1}
          \end{equation*}
          得到的线性方程组的解给出 \( r_1=2 \) 与 \( r_2=1 \)。
        \partsitem 是；由
          \( r_1(x^2)+r_2(2x+x^2)+r_3(x+x^3)=x-x^3 \)
          \begin{equation*}
            \begin{linsys}{3}
                  &  &2r_2 &+ &r_3 &= &1  \\
              r_1 &+ &r_2  &  &    &= &0  \\
                  &  &     &  &r_3 &= &-1
            \end{linsys}
          \end{equation*}
          可得 \( -1(x^2)+1(2x+x^2)-1(x+x^3)=x-x^3 \)。
       \partsitem 不是；所给两个矩阵的任意组合，其右上角元素都为零。
      \end{exparts}
    \end{answer}
  \item % \cite{Cleary} credit removed at request Cleary 2022-Sep-13
    一位超级英雄位于二维平面上的原点处。
    超级英雄有两件装置：一块悬浮滑板，可沿方向 $\binom{3}{1}$
    移动任意距离；
    以及一张魔毯，可沿方向
    $\binom{1}{2}$ 移动任意距离。
    \begin{exparts}
      \partsitem 一个邪恶反派藏身于平面上点
        $(-5,7)$ 处。
        超级英雄要移动多少个悬浮滑板单位和多少个魔毯单位，
        才能到达反派所在处？
     \partsitem 平面上是否存在某个地方，可以让反派安全藏身而不被找到？
        如果存在，给出一个这样的位置。
        如果不存在，解释原因。
     \partsitem 超级英雄和反派被传送到了一个三维空间，此时超级英雄拥有三件装置。
        \begin{equation*}
          \text{悬浮滑板: }\colvec{-1 \\ 0 \\ 3}
          \quad
          \text{魔毯: }\colvec{2 \\ 1 \\ 0}
          \quad
          \text{滑板车: }\colvec{5 \\ 4 \\ -9}
        \end{equation*}
        是否还有可以让反派安全藏身的地方？
        如果有，给出一个这样的位置；如果没有，解释原因。
    \end{exparts}
    \begin{answer}
      \begin{exparts}
        \partsitem 关系式
          \begin{equation*}
            \colvec{-5 \\ 7}=a\colvec{3 \\ 1}+b\colvec{1 \\ 2}
          \end{equation*}
          给出如下方程组。
          \begin{equation*}
            \begin{linsys}{2}
              3a &+ &b  &= &-5 \\
               a &+ &2b &= &7
            \end{linsys}
            \grstep{-(1/3)\rho_1+\rho_2}
            \begin{linsys}{2}
              3a &+ &b      &= &-5 \\
                 &  &(5/3)b &= &26/3
            \end{linsys}
          \end{equation*}
          所以 $b=26/5$。
          代入 $3a+(26/5)=-5$ 得 $a=-17/5$。
          用本小节的术语来说，向量 $\binom{-5}{7}$
          在另外两个向量构成的集合的张成之中。
        \partsitem 没有藏身之处。
          下面的高斯消元法化简
          \begin{equation*}
            \begin{linsys}{2}
              3a &+ &b  &= &x \\
               a &+ &2b &= &y
            \end{linsys}
            \grstep{-(1/3)\rho_1+\rho_2}
            \begin{linsys}{2}
              3a &+ &b      &= &x\hfill\hbox{} \\
                 &  &(5/3)b &= &-(1/3)x+y
            \end{linsys}
          \end{equation*}
          表明，对任意 $x,y\in\R$ 都存在一对 $a,b\in\R$。
          用本小节的术语来说，这两个向量的张成
          是整个平面。
        \partsitem 没有藏身之处。
          考虑这个方程组。
          \begin{equation*}
            a\colvec{-1 \\ 0 \\ 3}+b\colvec{2 \\ 1 \\ 0}+c\colvec{5 \\ 4 \\ -9}
                = \colvec{x \\ y \\ z}
          \end{equation*}
          化简后
          \begin{align*}
            \begin{linsys}{3}
              -a &+ &2b &+ &5c &= &x \\
                 &  &b  &+ &4c &= &y \\
              3a &  &   &- &9c &= &z \\
            \end{linsys}
            & \grstep{3\rho_1+\rho_3}
            \begin{linsys}{3}
              -a &+ &2b  &+ &5c &= &x\hfill\hbox{} \\
                 &  &b   &+ &4c &= &y\hfill\hbox{} \\
                 &  &6b  &+ &6c &= &3x+z \\
            \end{linsys}                                  \\
            & \grstep{-6\rho_2+\rho_3}
            \begin{linsys}{3}
              -a &+ &2b  &+ &5c  &= &x\hfill\hbox{} \\
                 &  &b   &+ &4c  &= &y\hfill\hbox{} \\
                 &  &    & &-18c &= &3x+-6y+z \\
            \end{linsys}
          \end{align*}
          表明，对任意 $x,y,z\in\R$，
          所需的 $a$、$b$ 与~$c$ 都存在。
          超级英雄的装置张成了整个空间。
      \end{exparts}
    \end{answer}
  \item
    下列哪些是实变量实值函数构成的向量空间中
    张成
    \( \spanof{\set{\cos^2x,\sin^2x} } \)
    的成员？
    \begin{exparts*}
      \partsitem \( f(x)=1 \)
      \partsitem \( f(x)=3+x^2 \)
      \partsitem \( f(x)=\sin x \)
      \partsitem \( f(x)=\cos (2x) \)
    \end{exparts*}
    \begin{answer}
      \begin{exparts}
        \partsitem 是；它在那个张成之中，因为
          \( 1\cdot\cos^2x+1\cdot\sin^2x=f(x) \)。
        \partsitem 不是，因为 \( r_1\cos^2x+r_2\sin^2x=3+x^2 \) 没有对
          所有 \( x \) 都成立的标量解。
          例如，取 $x$ 为 $0$ 和 $\pi$ 得到两个方程 $r_1\cdot 1+r_2\cdot 0=3$ 与
          $r_1\cdot 1+r_2\cdot 0=3+\pi^2$，二者
          相互矛盾。
        \partsitem 不是；考虑取 $x$ 为 $\pi/2$ 与
          $3\pi/2$ 时的结果。
        \partsitem 是，\( \cos (2x)=1\cdot\cos^2(x)-1\cdot\sin^2(x) \)。
      \end{exparts}
     \end{answer}
  \recommended \item
    下列哪些集合张成 \( \Re^3 \)？
    也就是说，下列哪些集合具有这样的性质：任意三元
    向量都可以表示为该集合元素的某个适当
    线性组合？
    \begin{exparts*}
      \partsitem \( \set{ \colvec[r]{1 \\ 0 \\ 0},
               \colvec[r]{0 \\ 2 \\ 0},
               \colvec[r]{0 \\ 0 \\ 3}  } \)
      \partsitem \( \set{ \colvec[r]{2 \\ 0 \\ 1},
               \colvec[r]{1 \\ 1 \\ 0},
               \colvec[r]{0 \\ 0 \\ 1}  } \)
      \partsitem \( \set{ \colvec[r]{1 \\ 1 \\ 0},
               \colvec[r]{3 \\ 0 \\ 0}  } \)
      \partsitem \( \set{ \colvec[r]{1 \\ 0 \\ 1},
               \colvec[r]{3 \\ 1 \\ 0},
               \colvec[r]{-1 \\ 0 \\ 0},
               \colvec[r]{2 \\ 1 \\ 5}  } \)
      \partsitem \( \set{ \colvec[r]{2 \\ 1 \\ 1},
               \colvec[r]{3 \\ 0 \\ 1},
               \colvec[r]{5 \\ 1 \\ 2},
               \colvec[r]{6 \\ 0 \\ 2}  } \)
    \end{exparts*}
    \begin{answer}
      \begin{exparts}
         \partsitem 是，对任意 \( x,y,z\in\Re \)，方程
          \begin{equation*}
             r_1\colvec[r]{1 \\ 0 \\ 0}
             +r_2\colvec[r]{0 \\ 2 \\ 0}
             +r_3\colvec[r]{0 \\ 0 \\ 3}
             =\colvec{x \\ y \\ z}
          \end{equation*}
          有解 \( r_1=x \)，\( r_2=y/2 \) 与
          \( r_3=z/3 \)。
        \partsitem 是，方程
          \begin{equation*}
            r_1\colvec[r]{2 \\ 0 \\ 1}
            +r_2\colvec[r]{1 \\ 1 \\ 0}
            +r_3\colvec[r]{0 \\ 0 \\ 1}
            =\colvec{x \\ y \\ z}
          \end{equation*}
          给出如下
          \begin{equation*}
            \begin{linsys}{3}
              2r_1 &+  &r_2  &  &    &=  &x \\
                   &   &r_2  &  &    &=  &y \\
               r_1 &   &     &+ &r_3 &=  &z \\
            \end{linsys}
            \grstep{-(1/2)\rho_1+\rho_3}\repeatedgrstep{(1/2)\rho_2+\rho_3}
            \begin{linsys}{3}
              2r_1 &+  &r_2  &  &    &=  &x\hfill\hbox{} \\
                   &   &r_2  &  &    &=  &y\hfill\hbox{} \\
                   &   &     &  &r_3 &=  &-(1/2)x+(1/2)y+z \\
            \end{linsys}
          \end{equation*}
          于是，给定任意 $x$、$y$ 和 $z$，我们可以算出
          \( r_3=(-1/2)x+(1/2)y+z \)、\( r_2=y \) 与
          \( r_1=(1/2)x-(1/2)y \)。
        \partsitem 不是。
          特别地，我们无法得到向量
          \begin{equation*}
            \colvec[r]{0 \\ 0 \\ 1}
          \end{equation*}
          的线性组合，因为所给的两个
          向量的第三个分量都为零。
       \partsitem 是。
         方程
         \begin{equation*}
           r_1\colvec[r]{1 \\ 0 \\ 1}
           +r_2\colvec[r]{3 \\ 1 \\ 0}
           +r_3\colvec[r]{-1\\ 0 \\ 0}
           +r_4\colvec[r]{2 \\ 1 \\ 5}
           =\colvec{x \\ y \\ z}
         \end{equation*}
         导出如下化简。
         \begin{equation*}
           \begin{amat}{4}
             1  &3  &-1  &2  &x  \\
             0  &1  &0   &1  &y  \\
             1  &0  &0   &5  &z
           \end{amat}
           \grstep{-\rho_1+\rho_3}\repeatedgrstep{3\rho_2+\rho_3}
           \begin{amat}{4}
             1  &3  &-1  &2  &x \\
             0  &1  &0   &1  &y  \\
             0  &0  &1   &6  &-x+3y+z
           \end{amat}
         \end{equation*}
         它有无穷多解。
         例如，我们可以取 $r_4$ 为零，然后按通常的
         回代方法，用 $x$、$y$ 和 $z$ 解出
         $r_3$、$r_2$ 和 $r_1$。
       \partsitem 不是。
         方程
         \begin{equation*}
           r_1\colvec[r]{2 \\ 1 \\ 1}
           +r_2\colvec[r]{3 \\ 0 \\ 1}
           +r_3\colvec[r]{5 \\ 1 \\ 2}
           +r_4\colvec[r]{6 \\ 0 \\ 2}
           =\colvec{x \\ y \\ z}
         \end{equation*}
         导出如下化简。
         \begin{multline*}
           \begin{amat}{4}
             2  &3  &5   &6  &x  \\
             1  &0  &1   &0  &y  \\
             1  &1  &2   &2  &z
           \end{amat}                                             \\
           \grstep[-(1/2)\rho_1+\rho_3]{-(1/2)\rho_1+\rho_2}
           \repeatedgrstep{-(1/3)\rho_2+\rho_3}
           \begin{amat}{4}
             2  &3     &5     &6  &x \\
             0  &-3/2  &-3/2  &-3 &-(1/2)x+y  \\
             0  &0     &0     &0  &-(1/3)x-(1/3)y+z
           \end{amat}
         \end{multline*}
         这表明并非每个三元向量都能如此表示。
         只有满足约束 $-(1/3)x-(1/3)y+z=0$ 的向量才在张成之中。
         （要看出任何这样的向量确实可以表示，
         可取 $r_3$ 与 $r_4$
         为零，然后用回代法以 $x$、$y$ 和
         $z$ 解出 $r_1$ 与 $r_2$。）
      \end{exparts}
    \end{answer}
  \recommended \item
    对每个子空间的描述给出参数化。
    然后把每个子空间表示为一个张成。
    \begin{exparts}
      \partsitem  三宽行向量中的子集
        \( \set{\rowvec{a &b &c}\suchthat a-c=0}   \)
      \partsitem \( \matspace_{\nbyn{2}} \) 的这个子集
        \begin{equation*}
          \set{\begin{mat}
                 a  &b  \\
                 c  &d
               \end{mat}  \suchthat a+d=0}
        \end{equation*}
      \partsitem \( \matspace_{\nbyn{2}} \) 的这个子集
        \begin{equation*}
          \set{\begin{mat}
                 a  &b  \\
                 c  &d
               \end{mat}  \suchthat \text{\( 2a-c-d=0 \)
                                               且 \( a+3b=0 \)} }
        \end{equation*}
      \partsitem 子集 \( \set{a+bx+cx^3\suchthat a-2b+c=0} \)，
        属于 \( \polyspace_3 \)
      \partsitem \( \polyspace_2 \) 的子集：满足 \( p(7)=0 \) 的
        次数不超过二的多项式 \( p \)
    \end{exparts}
    \begin{answer}
      \begin{exparts}
        \partsitem \( \set{\rowvec{c &b &c}\suchthat b,c\in\Re}
                      =\set{b\rowvec{0 &1 &0}+c\rowvec{1 &0 &1}
                        \suchthat b,c\in\Re} \)
           张成集合的显然选择是
           $\set{\rowvec{0 &1 &0},\rowvec{1 &0 &1}}$。
        \partsitem \( \set{\begin{mat}
                       -d &b  \\
                       c  &d
                      \end{mat} \suchthat b,c,d\in\Re}
                     =\set{b\begin{mat}[r]
                       0  &1  \\
                       0  &0
                      \end{mat}
                     +c\begin{mat}[r]
                       0  &0  \\
                       1  &0
                      \end{mat}
                     +d\begin{mat}[r]
                       -1  &0  \\
                       0  &1
                      \end{mat}  \suchthat b,c,d\in\Re} \)
            张成这个空间的一个集合就由这三个矩阵组成。
        \partsitem 方程组
          \begin{equation*}
            \begin{linsys}{4}
              a  &+  &3b  &   &   &  &  &=  &0  \\
             2a  &   &    &   &-c &- &d &=  &0
            \end{linsys}
          \end{equation*}
          给出 \( b=-(c+d)/6 \) 与 \( a=(c+d)/2 \)。
          于是一种描述如下。
          \begin{equation*}
            \set{c\begin{mat}[r]
                       1/2  &-1/6  \\
                       1    &0
                      \end{mat}
                     +d\begin{mat}[r]
                       1/2  &-1/6  \\
                       0    &1
                      \end{mat}  \suchthat c,d\in\Re}
           \end{equation*}
          这表明张成该子空间的一个集合由那
          两个矩阵组成。
        \partsitem 由 $a=2b-c$ 可知，集合
           \( \set{(2b-c)+bx+cx^3 \suchthat b,c\in\Re} \)
           等于集合
           \( \set{b(2+x)+c(-1+x^3) \suchthat b,c\in\Re}  \)。
           所以该子空间是集合 $\set{2+x, -1+x^3}$ 的张成。
        \partsitem 集合
          \( \set{a+bx+cx^2\suchthat a+7b+49c=0} \)
          可以参数化为
          \begin{equation*}
            \set{b(-7+x)+c(-49+x^2)\suchthat b,c\in\Re}
          \end{equation*}
          因此
          有张成集 $\set{-7+x,-49+x^2}$。
      \end{exparts}
    \end{answer}
  \recommended \item
    对给定空间中的给定子空间，求一个张成它的集合。
    （\textit{提示。}先对每一个作参数化。）
    \begin{exparts}
      \partsitem  \( \Re^3 \) 中的 \( xz \) 平面
      \partsitem \( \set{\colvec{x \\ y \\ z}\suchthat 3x+2y+z=0} \)
            在 \( \Re^3 \) 中
      \partsitem \( \set{\colvec{x \\ y \\ z \\ w}\suchthat
                       2x+y+w=0 \text{\ 且\ } y+2z=0} \)
            在 \( \Re^4 \) 中
      \partsitem \( \set{a_0+a_1x+a_2x^2+a_3x^3\suchthat
                        a_0+a_1=0 \text{\ 且\ } a_2-a_3=0} \)
            在 \( \polyspace_3 \) 中
      \partsitem 空间 \( \polyspace_4 \) 中的集合 \( \polyspace_4 \)
      \partsitem \( \matspace_{\nbyn{2}} \) 中的 \( \matspace_{\nbyn{2}} \)
    \end{exparts}
    \begin{answer}
      \begin{exparts}
        \partsitem 我们可以这样参数化
          \begin{equation*}
             \set{\colvec{x \\ 0 \\ z}\suchthat x,z\in\Re}
             =\set{x\colvec[r]{1 \\ 0 \\ 0}
                  +z\colvec[r]{0 \\ 0 \\ 1}\suchthat x,z\in\Re}
          \end{equation*}
          由此得到张成集如下。
          \begin{equation*}
             \set{\colvec[r]{1 \\ 0 \\ 0},\colvec[r]{0 \\ 0 \\ 1}}
          \end{equation*}
        \item 这是参数化描述，以及相应的张成集。
          \begin{equation*}
             \set{y\colvec[r]{-2/3 \\ 1 \\ 0}+z\colvec[r]{-1/3 \\ 0 \\ 1}
                 \suchthat y,z\in\Re }
             \qquad
             \set{\colvec[r]{-2/3 \\ 1 \\ 0},\colvec[r]{-1/3 \\ 0 \\ 1} }
          \end{equation*}
        \partsitem \( \set{\colvec[r]{1 \\ -2 \\ 1 \\ 0},
                      \colvec[r]{-1/2 \\ 0 \\ 0 \\ 1} } \)
        \partsitem 把描述参数化为
          \( \set{-a_1+a_1x+a_3x^2+a_3x^3\suchthat a_1,a_3\in\Re } \)
          得到 \( \set{-1+x,x^2+x^3}. \)
        \partsitem \( \set{1,x,x^2,x^3,x^4} \)
        \partsitem \( \set{ \begin{mat}[r]
                   1  &0  \\
                   0  &0
                 \end{mat},
                 \begin{mat}[r]
                   0  &1  \\
                   0  &0
                 \end{mat},
                 \begin{mat}[r]
                   0  &0  \\
                   1  &0
                 \end{mat},
                 \begin{mat}[r]
                   0  &0  \\
                   0  &1
                 \end{mat} } \)
      \end{exparts}
    \end{answer}
  \item
    \( \Re^2 \) 是 \( \Re^3 \) 的子空间吗？
    \begin{answer}
      严格地说，不是。
      \( \Re^3 \) 的子空间都是三元向量构成的集合，而
      \( \Re^2 \) 是二元向量构成的集合。
      不过显然，\( \Re^2 \) 与 \( \Re^3 \) 的这个子空间“几乎
      一样”。
      \begin{equation*}
        \set{\colvec{x \\ y \\ 0}\suchthat x,y\in\Re}
      \end{equation*}
    \end{answer}
  \recommended \item
    判断下列每一个是否为实变量实值函数构成的向量空间的子空间。
    \begin{exparts}
      \partsitem \definend{偶}\index{function!even}\index{even functions} 函数
        \( \set{\map{f}{\Re}{\Re} \suchthat f(-x)=f(x) \text{对一切 } x} \)。
        例如，该集合的两个成员是 $f_1(x)=x^2$
        与 $f_2(x)=\cos (x)$。
      \partsitem \definend{奇}\index{function!odd}\index{odd function}
        函数
        \( \set{\map{f}{\Re}{\Re} \suchthat f(-x)=-f(x) \text{对一切 } x} \)。
        两个成员是 $f_3(x)=x^3$ 与 $f_4(x)=\sin(x)$。
    \end{exparts}
    \begin{answer}
      当然，加法与标量乘法运算就是
      从外围空间继承的那两个运算。
      \begin{exparts}
        \partsitem 这是一个子空间。
          它不是空集，因为它至少包含上面给出的两个示例函数。
          它是封闭的，因为若 \( f_1,f_2 \) 是偶函数且
          \( c_1,c_2 \) 是标量，则有下式。
          \begin{equation*}
            (c_1f_1+c_2f_2)\,(-x)
            =c_1\,f_1(-x)+c_2\,f_2(-x)
            =c_1\,f_1(x)+c_2\,f_2(x)
            =(c_1f_1+c_2f_2)\,(x)
          \end{equation*}
        \partsitem 这也是一个子空间；验证方法与
          前一个类似。
      \end{exparts}
    \end{answer}
  \item
    \nearbyexample{ex:SpanSingVec} 说明，对向量空间 $V$ 中任一向量 $\vec{v}$，
    集合 $\set{r\cdot\vec{v}\suchthat r\in\Re}$
    是 $V$ 的一个子空间。
    （这不过是单元素集 $\set{\vec{v}}$ 的
    张成\index{span!of a singleton}。）
    这样的子空间是否一定是真子空间？
    \begin{answer}
      它可以不是真子空间。
      例如，
      若向量空间是 \( \Re^1 \) 且 \( \vec{v}\neq\zero\, \)，
      则子空间 $\set{r\cdot\vec{v}\suchthat r\in\Re}$
      就是整个 $\Re^1$。
    \end{answer}
  \item
    向量空间定义之后的一个例子表明，
    齐次线性方程组的解集是一个向量空间。
    用本小节的术语来说，当方程组有 $n$ 个变量时，它是 $\Re^n$ 的一个子空间。
    那么非齐次线性方程组呢；它的解是否构成一个
    子空间（在继承的运算下）？
    \begin{answer}
      不构成，这样的集合不封闭。
      首先，它不包含零向量。
    \end{answer}
  \item \cite{Cleary}
   就下列每一条给出一个例子，或者解释为什么
   不可能给出。
   \begin{exparts}
     \item $\matspace_{\nbyn{2}}$ 的一个不是
       子空间的非空子集。
     \item $\Re^2$ 中一个不张成该空间的两向量集合。
   \end{exparts}
   \begin{answer}
     \begin{exparts}
       \item $\matspace_{\nbyn{2}}$ 的这个非空子集不是子空间。
         \begin{equation*}
           A=\set{
             \begin{mat}[r]
               1 &2 \\
               3 &4
             \end{mat},
             \begin{mat}[r]
               5 &6 \\
               7 &8
             \end{mat}}
         \end{equation*}
         它不是 $\matspace_{\nbyn{2}}$ 的子空间的一个理由是
         它不包含零矩阵。
         （另一个理由是它对加法不封闭，因为两者的和
         不是 $A$ 的元素。
         它对标量乘法也不封闭。）
        \item 这个由两个向量构成的集合不张成 $\Re^2$。
          \begin{equation*}
            \set{\colvec[r]{1 \\ 1},\colvec[r]{3 \\ 3}}
          \end{equation*}
          这两个向量的任何线性组合都无法给出
          第二个分量不等于第一个分量的向量。
      \end{exparts}
   \end{answer}
  \item
    \nearbyexample{ex:SubspRThree} 表明
    $\Re^3$ 有无穷多个子空间。
    是否每个非平凡空间都有无穷多个子空间？
    \begin{answer}
      不是。
      \( \Re^1 \) 仅有的子空间是它本身与其
      平凡子空间。
      $\Re$ 的任何包含非零成员 $\vec{v}$ 的子空间 $S$
      必须包含其全部标量倍数构成的集合
      $\set{r\cdot\vec{v}\suchthat r\in\Re}$。
      但这个集合就是整个 $\Re$。
    \end{answer}
  \item \label{exer:SubspIffClosed}
    补全 \nearbylemma{th:SubspIffClosed} 的证明。
    \begin{answer}
      第~(1) 项已在正文中验证。

      第~(2) 条有五个条件。
      第一条是封闭性：若 \( c\in\Re \) 且 \( \vec{s}\in S \)，则
      \( c\cdot\vec{s}\in S \)，因为
      \( c\cdot\vec{s}=c\cdot\vec{s}+0\cdot\zero \)。
      第二，因为 \( S \) 中的运算继承自 \( V \)，所以对 \( c,d\in\Re \) 与 \( \vec{s}\in S \)，
      \( S \) 中的标量乘积 \( (c+d)\cdot\vec{s}\, \) 等于 \( V \) 中的乘积
      \( (c+d)\cdot\vec{s}\, \)，而它等于 \( V \) 中的
      \( c\cdot\vec{s}+d\cdot\vec{s}\, \)，进而等于 \( S \) 中的
      \( c\cdot\vec{s}+d\cdot\vec{s}\, \)。

      第三、第四和第五个条件的验证与刚才给出的第二个条件的验证类似。
    \end{answer}
  \item 
    证明每个向量空间只有一个平凡子空间。
    \begin{answer}
      前一小节的一道练习证明了每个向量空间只有一个零向量（也就是说，只有
      一个向量是该空间的加法单位元）。
      而平凡空间只有一个元素，且该元素必是这个（唯一的）零向量。
    \end{answer}
  \item 
    证明对向量空间的任意子集 \( S \)，张成的张成等于张成
    \( \spanof{ \spanof{S} }=\spanof{S} \)。
    （\textit{提示。}
    $\spanof{S}$ 的成员是 $S$ 的成员的线性组合。
    $\spanof{\spanof{S}}$ 的成员是 $S$ 的成员的线性组合的
    线性组合。）
    \begin{answer}
      如提示所言，根本的原因是第一章的线性组合引理。
      为给出完整证明，我们将证明这两个集合互相包含。

      第一个包含关系 \( \spanof{ \spanof{S} }\supseteq\spanof{S}  \)
      是一个更一般且显然的事实的特例：对向量空间的任意子集
      \( T \)，都有 \( \spanof{T}\supseteq T \)。

      对于另一个包含关系，
      即 \( \spanof{ \spanof{S} }\subseteq\spanof{S} \)，
      从 \( \spanof{S} \) 中取 $m$ 个向量，即
      \( c_{1,1}\vec{s}_{1,1}+\cdots+c_{1,n_1}\vec{s}_{1,n_1} \)、\ldots、
      \( c_{1,m}\vec{s}_{1,m}+\cdots+c_{1,n_m}\vec{s}_{1,n_m} \)，
      并注意它们的任意线性组合
      \begin{equation*}
         r_1(c_{1,1}\vec{s}_{1,1}+\cdots+c_{1,n_1}\vec{s}_{1,n_1})+\cdots
        +r_m(c_{1,m}\vec{s}_{1,m}+\cdots+c_{1,n_m}\vec{s}_{1,n_m})
      \end{equation*}
      是 \( S \) 中元素的线性组合
      \begin{equation*}
        = (r_1c_{1,1})\vec{s}_{1,1}+\cdots+(r_1c_{1,n_1})\vec{s}_{1,n_1}+\cdots
        +(r_mc_{1,m})\vec{s}_{1,m}+\cdots+(r_mc_{1,n_m})\vec{s}_{1,n_m}
      \end{equation*}
      因而在 \( \spanof{S} \) 中。
      也就是说，只需回想：线性组合的线性组合（$S$ 中成员的）
      仍是线性组合（同样还是 $S$ 中成员的）。
    \end{answer}
  \item 
    我们见过的所有子空间都在其描述中以某种方式用到了零。
    例如，\nearbyexample{ex:SubspacesRTwo} 中的子空间由 $\Re^2$ 中
    第二个分量为零的所有向量组成。
    相比之下，
    $\Re^2$ 中第二个分量为 1 的向量集合
    不构成子空间（它对标量乘法不封闭）。
    另一个例子是 \nearbyexample{ex:PlaneSubspRThree}，其中对向量的条件
    是三个分量之和为零。
    如果那里的条件
    改为三个分量之和为 1，
    那么它就不是子空间
    （同样，它不封闭）。
    不过，对零的依赖并非严格必需。
    考虑集合
    \begin{equation*}
       \set{\colvec{x \\ y \\ z}\suchthat x+y+z=1}
    \end{equation*}
    在这些运算下。
    \begin{equation*}
       \colvec{x_1 \\ y_1 \\ z_1}+\colvec{x_2 \\ y_2 \\ z_2}
       =\colvec{x_1+x_2-1 \\ y_1+y_2 \\ z_1+z_2}
       \qquad
       r\colvec{x \\ y \\ z}=\colvec{rx-r+1 \\ ry \\ rz}
    \end{equation*}
    \begin{exparts}
      \partsitem 证明它不是 $\Re^3$ 的子空间。
         （\textit{提示。}   参见 \nearbyexample{ex:OperNotInherit}）。
      \partsitem 证明它是一个向量空间。         
         注意由前一问，
         \nearbylemma{th:SubspIffClosed} 不能适用。  
      \partsitem 证明 $\Re^3$ 的任何子空间必定经过原点，
         因此 $\Re^3$ 的任何子空间在其描述中必定涉及零。
         逆命题成立吗？  
          $\Re^3$ 的任何包含原点的子集，
          在给定继承的运算后会成为子空间吗？
    \end{exparts}
    \begin{answer}
      \begin{exparts}
         \partsitem 它不是子空间，因为这些运算不是继承来的运算。  
           例如，在这个空间中，
           \begin{equation*}
             0\cdot\colvec{x \\ y \\ z}=\colvec[r]{1 \\ 0 \\ 0}
           \end{equation*}
           而这当然在 $\Re^3$ 中不成立。
         \partsitem 我们可以把证明对加法封闭的论证
           与证明对标量乘法封闭的论证合并成
           一个证明对两个向量的线性组合封闭的单一论证。
           若 $r_1,r_2,x_1,x_2,y_1,y_2,z_1,z_2$ 都属于 $\Re$，则      
           \begin{multline*}
              r_1\colvec{x_1 \\ y_1 \\ z_1}
              +r_2\colvec{x_2 \\ y_2 \\ z_2}
              =\colvec{r_1x_1-r_1+1 \\ r_1y_1 \\ r_1z_1}
               +\colvec{r_2x_2-r_2+1 \\ r_2y_2 \\ r_2z_2}                 \\
            =\colvec{r_1x_1-r_1+r_2x_2-r_2+1 \\ r_1y_1+r_2y_2 \\ r_1z_1+r_2z_2}
           \end{multline*} 
           （注意这个空间中加法的定义是：第一个分量按
           $(r_1x_1-r_1+1)+(r_2x_2-r_2+1)-1$ 的方式合并，
           所以最后一个向量的第一个分量并非表示
           `$\hbox{}+2$'）。
           将最后一个向量的三个分量相加得
           $r_1(x_1-1+y_1+z_1)+r_2(x_2-1+y_2+z_2)+1=r_1\cdot0+r_2\cdot0+1=1$。

           其余大部分条件的验证都很容易（尽管这些运算的古怪之处
           使它们无法按常规进行）。
           加法的交换律如下。
           \begin{equation*}
             \colvec{x_1 \\ y_1 \\ z_1}+\colvec{x_2 \\ y_2 \\ z_2}
             =\colvec{x_1+x_2-1 \\ y_1+y_2 \\ z_1+z_2}
             =\colvec{x_2+x_1-1 \\ y_2+y_1 \\ z_2+z_1}
             =\colvec{x_2 \\ y_2 \\ z_2}+\colvec{x_1 \\ y_1 \\ z_1}
           \end{equation*}
           加法的结合律则是
           \begin{equation*}
             (\colvec{x_1 \\ y_1 \\ z_1}+\colvec{x_2 \\ y_2 \\ z_2})
              +\colvec{x_3 \\ y_3 \\ z_3}
             =\colvec{(x_1+x_2-1)+x_3-1 \\ (y_1+y_2)+y_3 \\ (z_1+z_2)+z_3}
           \end{equation*}
           而
           \begin{equation*}
             \colvec{x_1 \\ y_1 \\ z_1}
             +(\colvec{x_2 \\ y_2 \\ z_2}+\colvec{x_3 \\ y_3 \\ z_3})
             =\colvec{x_1+(x_2+x_3-1)-1 \\ y_1+(y_2+y_3) \\ z_1+(z_2+z_3)}
           \end{equation*}
           且两者相等。
           关于这个加法运算的单位元 
           如下
           \begin{equation*}
             \colvec{x \\ y \\ z}+\colvec[r]{1 \\ 0 \\ 0}
             =\colvec{x+1-1 \\ y+0 \\ z+0}
             =\colvec{x \\ y \\ z}
           \end{equation*}
           加法逆元类似。
           \begin{equation*}
             \colvec{x \\ y \\ z}+\colvec{-x+2 \\ -y \\ -z}
             =\colvec{x+(-x+2)-1 \\ y-y \\ z-z}
             =\colvec[r]{1 \\ 0 \\ 0}
           \end{equation*}

           关于标量乘法的那些条件也容易验证。
           对第一个条件，
           \begin{equation*}
             (r+s)\colvec{x \\ y \\ z}
             =\colvec{(r+s)x-(r+s)+1 \\ (r+s)y \\ (r+s)z}
           \end{equation*}
           而
           \begin{multline*}
             r\colvec{x \\ y \\ z}+s\colvec{x \\ y \\ z}
             =\colvec{rx-r+1 \\ ry \\ rz}+\colvec{sx-s+1 \\ sy \\ sz}  \\
             =\colvec{(rx-r+1)+(sx-s+1)-1 \\ ry+sy \\ rz+sz}
           \end{multline*}
           两者相等。
           第二个条件比较
           \begin{equation*}
             r\cdot(\colvec{x_1 \\ y_1 \\ z_1}+\colvec{x_2 \\ y_2 \\ z_2})
             =r\cdot\colvec{x_1+x_2-1 \\ y_1+y_2 \\ z_1+z_2}
             =\colvec{r(x_1+x_2-1)-r+1 \\ r(y_1+y_2) \\ r(z_1+z_2)}
           \end{equation*}
           与
           \begin{multline*}
             r\colvec{x_1 \\ y_1 \\ z_1}+r\colvec{x_2 \\ y_2 \\ z_2}
             =\colvec{rx_1-r+1 \\ ry_1 \\ rz_1}
                     +\colvec{rx_2-r+1 \\ ry_2 \\ rz_2}                  \\
             =\colvec{(rx_1-r+1)+(rx_2-r+1)-1 \\ ry_1+ry_2 \\ rz_1+rz_2}
           \end{multline*}
           且它们相等。
           对第三个条件，
           \begin{equation*}
             (rs)\colvec{x \\ y \\ z}
             =\colvec{rsx-rs+1 \\ rsy \\ rsz}
           \end{equation*}
           而
           \begin{equation*}
             r(s\colvec{x \\ y \\ z})
             =r(\colvec{sx-s+1 \\ sy \\ sz})
             =\colvec{r(sx-s+1)-r+1 \\ rsy \\ rsz}
           \end{equation*}
           两者相等。
           对于被 $1$ 的标量乘法，我们有下式。
           \begin{equation*}
             1\cdot\colvec{x \\ y \\ z}
             =\colvec{1x-1+1 \\ 1y \\ 1z}
             =\colvec{x \\ y \\ z}
           \end{equation*}
           这样，这两个运算满足向量空间的所有条件。

           \textit{注记。}
           理解这个向量空间的一种方式是把它看作 $\Re^3$ 中的平面
           \begin{equation*}
             P=\set{\colvec{x \\ y \\ z}\suchthat x+y+z=0}
           \end{equation*}
           沿 $x$ 轴离开原点平移 $1$ 所得的平面。
           于是加法变成：~要加上这个空间的两个成员
           \begin{equation*}
             \colvec{x_1 \\ y_1 \\ z_1},\;\colvec{x_2 \\ y_2 \\ z_2}
           \end{equation*}
           （其中 $x_1+y_1+z_1=1$ 且 $x_2+y_2+z_2=1$）
           先把它们移回 $1$ 放入 $P$，再按通常方式相加，
           \begin{equation*}
             \colvec{x_1-1 \\ y_1 \\ z_1}+\colvec{x_2-1 \\ y_2 \\ z_2}
             =\colvec{x_1+x_2-2 \\ y_1+y_2 \\ z_1+z_2}
             \qquad\text{（在 $P$ 中）}
           \end{equation*}
           然后把结果沿 $x$ 轴再移出 $1$。
           \begin{equation*}
             \colvec{x_1+x_2-1 \\ y_1+y_2 \\ z_1+z_2}.
           \end{equation*}
           标量乘法类似。
         \partsitem 为使该子空间对继承的标量
           乘法封闭，设 $\vec{v}$ 是该子空间的一个成员，
           \begin{equation*}
             0\cdot\vec{v}=\colvec[r]{0 \\ 0 \\ 0}
           \end{equation*}
           它也必须是成员。

           逆命题不成立。
           下面是 $\Re^3$ 的一个包含原点的子集
           \begin{equation*}
             \set{\colvec[r]{0 \\ 0 \\ 0},\colvec[r]{1 \\ 0 \\ 0}}
           \end{equation*}
           （这个子集只有两个元素）但它不是子空间。
       \end{exparts}
     \end{answer}
  \item 
    我们可以为“零个向量之和等于零向量”这一约定给出一个理由。
    考虑三个向量的和 $\vec{v}_1+\vec{v}_2+\vec{v}_3$。
    \begin{exparts}
      \partsitem 这个三个向量的和与这三个中前两个的和之差是什么？
      \partsitem 前一个和与仅第一个向量的和之差是什么？
      \partsitem 前一个只含一个向量的和与零个向量的和之差应当是什么？
      \partsitem 那么零个向量的和应当定义为什么？
    \end{exparts}
    \begin{answer}
      \begin{exparts}
        \item $(\vec{v}_1+\vec{v}_2+\vec{v}_3)-(\vec{v}_1+\vec{v}_2)
                 =\vec{v}_3$
        \item $(\vec{v}_1+\vec{v}_2)-(\vec{v}_1)
                 =\vec{v}_2$
        \item 显然，$\vec{v}_1$。
        \item 取那个只含一项的和再相减，得
          （$\vec{v}_1)-\vec{v}_1=\zero$。
      \end{exparts}
    \end{answer}
  \item 
    一个空间由它的子空间决定吗？
    也就是说，如果两个向量空间有相同的子空间，这两个空间是否必定相等？
    \begin{answer}
      是的；任何空间都是它自身的子空间，所以每个空间都包含另一个。  
    \end{answer}
  \item 
     \begin{exparts}
      \partsitem 给出一个对标量乘法封闭
        但对加法不封闭的集合。
      \partsitem 给出一个对加法封闭但对标量乘法
        不封闭的集合。
      \partsitem 给出一个对两者都不封闭的集合。
    \end{exparts}
    \begin{answer}
      \begin{exparts}
        \partsitem \( \Re^2 \) 中 \( x \) 轴与 \( y \) 轴的并就是一例。
        \partsitem 整数集，作为 \( \Re^1 \) 的子集，就是一例。
        \partsitem \( \Re^2 \) 的子集 \( \set{\vec{v}} \) 就是一例，
          其中 $\vec{v}$ 是任意非零向量。
      \end{exparts}  
     \end{answer}
  \item 
    证明一组向量的张成不依赖于列举这些向量时的顺序，即
    该集合中向量被列举的顺序。
    \begin{answer}
      由于向量空间加法是交换的，重新排列被加项不会改变线性组合。  
    \end{answer}
  \item  
    空集的张成是哪个平凡子空间？
    是
    \begin{equation*}
      \set{\colvec[r]{0 \\ 0 \\ 0}}\subseteq \Re^3,
      \quad\text{或}\quad
      \set{0+0x}\subseteq \polyspace_1,
    \end{equation*}
    还是某个别的子空间？
    \begin{answer}
      我们总是在一个包围空间的背景下考虑这个张成。  
    \end{answer}
  \item   
    证明若一个向量在某个集合的张成中，则把该向量加进集合
    不会使张成变大。
    这在“仅当”方向上也成立吗？
    \begin{answer}
      “当”与“仅当”两个方向都成立。

      对于“当”的方向，
      设 \( S \) 是向量空间 \( V \) 的子集，并假设
      \( \vec{v}\in S \) 满足
      \( \vec{v}=c_1\vec{s}_1+\dots+c_n\vec{s}_n \)，其中
      \( c_1,\ldots,c_n \) 是标量，且
      \( \vec{s}_1,\ldots,\vec{s}_n\in S \)。
      我们必须证明 \( \spanof{S\union\set{\vec{v}} }=\spanof{S} \)。

      一侧的包含关系，
      \( \spanof{S}\subseteq\spanof{S\union\set{\vec{v}} } \) 是显然的。
      对于另一个方向，
      \( \spanof{S\union\set{\vec{v}} }\subseteq\spanof{S} \)，注意若一个
      向量属于左边的集合，则它形如
      \( d_0\vec{v}+d_1\vec{t}_1+\dots+d_m\vec{t}_m \)，其中诸 \( d \) 是
      标量，诸 \( \vec{t}\, \) 属于 \( S \)。
      把它改写为
      \( d_0(c_1\vec{s}_1+\dots+c_n\vec{s}_n)
      +d_1\vec{t}_1+\cdots+d_m\vec{t}_m \)，并注意
      该结果是 \( S \) 的张成中的一个成员。

      “仅当”显然成立\Dash 加入 \( \vec{v} \) 
      会把张成扩大到至少包含 \( \vec{v} \)。
    \end{answer}
  \recommended \item
    子空间是子集，所以我们自然会考虑“是……的子空间”这一关系
    与通常的集合运算如何相互作用。
    \begin{exparts}
      \partsitem 若 \( A,B \) 是向量空间的子空间，它们的交
        \( A\intersection B \) 必是子空间吗？
        总是？有时？从不？
      \partsitem 并 \( A\union B \) 必是子空间吗？
      \partsitem 若 \( A \) 是某个~$V$ 的子空间，它的补集 $V-A$
        必是子空间吗？
    \end{exparts}
    （\textit{提示。}   试用 \nearbyexample{ex:SubspRThree} 中
    的一些试验性子空间。）
    \begin{answer}
      \begin{exparts}
        \partsitem 总是。

          设 \( A,B \) 是 \( V \) 的子空间。
          注意 
          它们的交不空，因为两者都包含零向量。
          若 \( \vec{w},\vec{s}\in A\intersection B \) 且 \( r,s \) 是
          标量，则 \( r\vec{v}+s\vec{w}\in A \)，因为
          每个向量都属于 \( A \)，所以线性组合属于 \( A \)，
          出于同样的理由 \(r\vec{v}+s\vec{w}\in B \)。
          于是这个交是封闭的。
          这时 \nearbylemma{th:SubspIffClosed} 适用。
        \partsitem 有时（更准确地说，仅当 \( A\subseteq B \) 或
          \( B\subseteq A \)）。

          为看出答案不是“总是”，取 \( V \) 为 \( \Re^3 \)，
          取 \( A \) 为 $x$ 轴，\( B \) 为
          \( y \) 轴。
          注意
          \begin{equation*}
            \colvec[r]{1 \\ 0}\in A \text{ 且 }\colvec[r]{0 \\ 1}\in B
            \quad\text{而}\quad
            \colvec[r]{1 \\ 0}+\colvec[r]{0 \\ 1}\not\in A\union B
          \end{equation*}
          因为该和既不在 \( A \) 中也不在 \( B \) 中。

          答案不是“从不”，因为若 \( A\subseteq B \) 或
          \( B\subseteq A \)，则显然 \( A\union B \) 是子空间。

          为证明 \( A\union B \) 是子空间仅当一个子空间
          包含另一个，我们假设 \( A\not\subseteq B \)
          且 \( B\not\subseteq A \)，并证明 
          该并不是子空间。
          \( A \) 不是 \( B \) 的子集这一假设意味着存在 
          \( \vec{a}\in A \) 且 \( \vec{a}\not\in B \)。
          另一假设给出一个 \( \vec{b}\in B \)，满足
          \( \vec{b}\not\in A \)。
          考虑 \( \vec{a}+\vec{b} \)。
          注意该和不是 \( A \) 的元素，否则 
          \( (\vec{a}+\vec{b})-\vec{a} \) 将属于 \( A \)，但它不属于。
          类似地，该和也不是 \( B \) 的元素。
          因此该和不是 \( A\union B \) 的元素，所以该并
          不是子空间。
        \partsitem 从不。
          由于 \( A \) 是子空间，它包含零向量，因此
          作为 $A$ 的补集的集合 $V-A$ 不包含它。
          没有零向量，这个补集不可能是向量空间。
      \end{exparts}  
    \end{answer}
  \item 
    一个集合的张成依赖于包围它的空间吗？
    也就是说，若 \( W \) 是 \( V \) 的子空间，\( S \) 是
    \( W \) 的子集（因而也是 \( V \) 的子集），那么 \( S \) 在
    \( W \) 中的张成可能不同于 \( S \) 在 \( V \) 中的张成吗？
    \begin{answer}
      一个集合的张成不依赖于包围它的空间。
      \( S \) 中向量的线性组合给出同样的和，
      无论我们把运算视为 \( W \) 的还是 \( V \) 的，
      因为 \( W \) 的运算继承自 \(  V \)。  
    \end{answer}
  \item 
    “是……的子空间”这一关系是传递的吗？
    也就是说，若 $V$ 是 $W$ 的子空间，$W$ 是 $X$ 的
    子空间，$V$ 必是 $X$ 的子空间吗？ 
    \begin{answer}
      是传递的；
      应用 \nearbylemma{th:SubspIffClosed}。
      （你必须考虑下面这一点。
      设 \( B \) 是向量空间 \( V \) 的子空间，再设
      \( A\subseteq B\subseteq V \) 是一个子空间。
      \( A \) 从哪个空间继承它的运算？
      答案是这无关紧要\Dash 两种情形下 \( A \) 继承的
      都是相同的运算。）  
    \end{answer}
  \item
    由于“张成”是集合上的一种运算，我们自然会考虑
    它与通常的集合运算如何相互作用。
    \begin{exparts}
      \partsitem 若 \( S\subseteq T \) 是向量空间的子集，是否有
        \( \spanof{S}\subseteq\spanof{T} \)？
        总是？有时？从不？
      \partsitem 若 \( S,T \) 是向量空间的子集，是否有
        \( \spanof{S\union T}=\spanof{S}\union\spanof{T} \)？
      \partsitem 若 \( S,T \) 是向量空间的子集，是否有
        \( \spanof{S\intersection T}=\spanof{S}\intersection\spanof{T} \)？
      \partsitem 补集的张成等于张成的补吗？
    \end{exparts}
    \begin{answer}
      \begin{exparts}
         \partsitem 总是；
           若 \( S\subseteq T \)，则 \( S \) 中元素的线性组合
           也是 \( T \) 中元素的线性组合。
         \partsitem 有时（更准确地说，当且仅当 
           \( S\subseteq T \) 或 \( T\subseteq S \)）。

           答案不是“总是”，\( \Re^3 \) 中的这个例子
           表明了这一点
           \begin{equation*}
             S=\set{\colvec[r]{1 \\ 0 \\ 0},\colvec[r]{0 \\ 1 \\ 0}},\quad
             T=\set{\colvec[r]{1 \\ 0 \\ 0},\colvec[r]{0 \\ 0 \\ 1}}
           \end{equation*}
           因为有下式。
           \begin{equation*}
             \colvec[r]{1 \\ 1 \\ 1}\in\spanof{S\union T}
             \qquad
             \colvec[r]{1 \\ 1 \\ 1}\not\in\spanof{S}\union \spanof{T}
           \end{equation*}

           答案不是“从不”，因为若其中一个集合包含另一个，
           则等式显然成立。
           我们可以通过假设 \( S\not\subseteq T \)（意味着存在
           向量 \( \vec{s}\in S \) 满足 
           \( \vec{s}\not\in T \)）
           且 \( T\not\subseteq S \)（给出一个 \( \vec{t}\in T \)，满足
           \( \vec{t}\not\in S \)），并注意
           \( \vec{s}+\vec{t}\in\spanof{S\union T} \)，
           再证明 
           \( \vec{s}+\vec{t}\not\in\spanof{S}\union\spanof{T} \)，
           来刻画等式仅在其中一个集合包含另一个时才发生。
         \partsitem 有时。

           显然
           \( \spanof{S\intersection T}
             \subseteq\spanof{S}\intersection\spanof{T} \)
           因为来自 
           \( S\intersection T \)
           的向量的任何线性组合既是 \( S \) 中向量的组合，也是
           \( T \) 中向量的组合。

           反方向的包含并不总成立。
           例如，在 \( \Re^2 \) 中，取
           \begin{equation*}
             S=\set{\colvec[r]{1 \\ 0},\colvec[r]{0 \\ 1}},\quad
             T=\set{\colvec[r]{2 \\ 0}}
           \end{equation*}
           于是 \( \spanof{S}\intersection\spanof{T} \) 是 \( x \) 轴，
           但 \( \spanof{S\intersection T}  \) 是平凡子空间。

           要精确刻画等式何时成立并不容易。
           显然若其中一个集合包含另一个则等式成立，但由 \( \Re^3 \) 中的这个例子
           可知那不是“仅当”的。
           \begin{equation*}
             S=\set{\colvec[r]{1 \\ 0 \\ 0},\colvec[r]{0 \\ 1 \\ 0}},
             \quad
             T=\set{\colvec[r]{1 \\ 0 \\ 0},\colvec[r]{0 \\ 0 \\ 1}}
           \end{equation*}
         \partsitem 从不，因为补集的张成是子空间，而
           张成的补不是（它不包含零 
           向量）。
      \end{exparts}  
     \end{answer}
  % \item 
  %   Reprove \nearbylemma{le:SpanIsASubsp} without doing the
  %   empty set separately.
  %   \begin{answer}
  %     Call the subset \( S \).
  %     By \nearbylemma{th:SubspIffClosed},
  %     we need to check that 
  %     \( \spanof{S} \) is closed under linear combinations.
  %     If \( c_1\vec{s}_1+\dots+c_n\vec{s}_n,
  %       c_{n+1}\vec{s}_{n+1}+\dots+c_m\vec{s}_m\in\spanof{S} \) then
  %     for any \( p,r\in\Re \) we have
  %     \begin{multline*}
  %       p\cdot(c_1\vec{s}_1+\cdots+c_n\vec{s}_n)+
  %            r\cdot(c_{n+1}\vec{s}_{n+1}+\cdots+c_m\vec{s}_m)        \\
  %       =
  %       pc_1\vec{s}_1+\cdots+pc_n\vec{s}_n
  %         +rc_{n+1}\vec{s}_{n+1}+\cdots+rc_m\vec{s}_m
  %     \end{multline*}
  %     which is an element of \( \spanof{S} \).
  %     % (\textit{Remark.}
  %     % If the set $S$ is empty, then that
  %     % `if \ldots\ then \ldots' statement is vacuously true.)  
  %   \end{answer}
  \item 找一个对线性组合封闭 
    然而却不是向量空间的结构。
    % (\textit{Remark.} 
    % This is a bit of a trick question.)
    \begin{answer}
      要出现这种情况，保证加法与标量乘法运算合理性的条件中
      必须有一个被违反。
      考虑带这些运算的 \( \Re^2 \)。
      \begin{equation*}
        \colvec{x_1 \\ y_1}
        +\colvec{x_2 \\ y_2}
        =\colvec[r]{0 \\ 0}
        \qquad
        r\colvec{x \\ y}
        =\colvec[r]{0 \\ 0}
      \end{equation*}
      集合 $\Re^2$ 对这些运算封闭。 
      但它不是向量空间。
      \begin{equation*}
        1\cdot\colvec[r]{1 \\ 1}
        \neq\colvec[r]{1 \\ 1}
      \end{equation*}  
    \end{answer}
\index{subspace|)}
\index{vector space|)}
\end{exercises}
